\documentclass[12pt]{article}
\usepackage{latexsym}
\usepackage{amsmath,amsfonts,amssymb,epsfig}
\usepackage{amsthm,amscd}
\usepackage{mathrsfs} 
%\usepackage{hyperref}

%  Set the margins
\setlength{\oddsidemargin}{0in}
\setlength{\evensidemargin}{0in}
\setlength{\topmargin}{-0.25in}
%\setlength{\headheight}{12pt}
%\setlength{\headsep}{25pt}
%\setlength{\footheight}{12pt}
%\setlength{\footskip}{75pt}
\setlength{\textheight}{8.5in}
\setlength{\textwidth}{6.5in}


\newtheorem{theorem}{Theorem}[section]
\newtheorem{definition}[theorem]{Definition}
\newtheorem{corollary}[theorem]{Corollary}
\newtheorem{proposition}[theorem]{Proposition}
\newtheorem{lemma}[theorem]{Lemma}
\newtheorem{conjecture}[theorem]{Conjecture}
\newtheorem{claim}[theorem]{Claim}
\newtheorem{rem}[theorem]{Remark}
\newenvironment{remark}{\begin{rem} \rm}{\end{rem}}
\newtheorem{ide}[theorem]{Idea}
\newenvironment{idea}{\begin{ide} \rm}{\end{ide}}
\newtheorem{exa}[theorem]{Example}
\newenvironment{example}{\begin{exa} \rm}{\end{exa}}
\newtheorem{ques}[theorem]{Question}
\newenvironment{question}{\begin{ques} \rm}{\end{ques}}
\newtheorem{pro}[theorem]{Problem}
\newenvironment{problem}{\begin{pro} \rm}{\end{pro}}
\newtheorem{spe}[theorem]{Speculation}
\newenvironment{specu}{\begin{spe} \rm}{\end{spe}}


\font\bbb=msbm10 scaled 1100

\newcommand{\Lie}{{\cal L}}
\newcommand{\R}{\mbox{\bbb R}}       % {\mbox{\rm I} \hspace{-.03in} {\bf R}}
\newcommand{\C}{\mbox{\bbb C}}       %{\mbox{\rm I} \hspace{-.09in} {\bf C}}
\newcommand{\Hp}{\mbox{\bbb H}}
\newcommand{\Z}{\mbox{\bbb Z}}
\newcommand{\Cnf}{\text{Conf}}

\newcommand{\Sp}{\mbox{\bbb S}}
\newcommand{\df}[1]{{\sc #1}}
\newcommand{\bfx}{{\mbox{\bf}}}
\newcommand{\no}{\noindent}


%%%%%%%%%%%%%%%%%%%%%%%%%%%%%%% begin %%%%%%%%%%%%%%%%%%%%%%%%%%%%%%%%%%%

\title{On the third order helicity of magnetic fields\\ on invariant domains in $S^3$} 

\author{R. Komendarczyk\footnote{Department of Mathematics,\ University of Pennsylvania; e-mail: \emph{rako@math.upenn.edu}.
 This project is supported  by DARPA, \#FA9550-08-1-0386.
}}

\begin{document}
\maketitle

\begin{abstract}
 \no We introduce an alternative approach to the third order helicity of a volume preserving vector field $X$. The proposed approach exploits correspondence between the Milnor $\bar{\mu}_{123}$-invariant for 3-component links and the homotopy invariants of maps to configuration spaces, and we provide a simple geometric proof of this fact in the case of borromean links. 
 Based on these connections we develop a formulation for the third order helicity of $X$ on invariant \emph{unlinked} domains of $X$, and 
 provide Arnold's style ergodic interpretation of this invariant
 as an average asymptotic $\bar{\mu}_{123}$-invariant of orbits of $X$. 
 We also obtain a lower bound for the $L^2$-energy of $X$.
\end{abstract} 

\section{Introduction}
 A purpose of this paper is to develop a particular formula for the third order helicity on certain  invariant sets of a volume preserving vector field $X$. The third order helicity, \cite{Khesin98}, is an invariant of $X$ under the action of volumorphisms isotopic to the identity (denoted here by $\text{Diff}^\text{vol}_0(M)$).  
\no Importance of such invariants stems from the basic fact that e.g. the evolution of the vorticity in the ideal hydrodynamics or the magnetic field $\mathbf{B}_0$ in the ideal magnetohydrodynamics (MHD), occurs along a path $t\longrightarrow g(t)\in \text{Diff}^\text{vol}_0(M)$, \cite[p. 176]{Khesin98}. Namely,
$\mathbf{B}(t)=g_\ast(t)\mathbf{B}_0$ which is a direct consequence of Euler's equations: 
\begin{equation}\label{eq:Eulers-eq}
\frac{d}{dt} \mathbf{B}+[\mathbf{v},\mathbf{B}]=0,\qquad \frac{d}{dt} g(t)=\mathbf{v} .
\end{equation}
 One often says that the magnetic field $\mathbf{B}$ is \emph{frozen in} the velocity field $\mathbf{v}$ of plasma, the action by $\text{Diff}^\text{vol}_0(M)$ is often referred to as \emph{frozen-in-field} deformations. The fundamental example of such invariant defined for a general class in $\text{Vect}^\text{vol}(M)$ is the \emph{helicity} $\mathsf{H}_{12}(X_1,X_2)$ defined for a pair of vector fields $X_1$ and $X_2$ on $M=S^3$ or a homology sphere. 
Helicity has been first introduced by Woltjer, in the context of magnetic fields, and is a measure of how orbits of $X_1$, and $X_2$ link with each other. This topological interpretation of helicity has been made precise by Arnold, who introduced (in  \cite{Arnold86}) the concept of the average asymptotic linking number of orbits for a pair volume preserving vector fields on $M$ and proved that it is equal to their helicity. 
The subject has been further developed 
in \cite{Berger90, Akhmetiev05, Hornig04, Laurence-Stredulinsky00b} (see \cite{Khesin98} for additional references), where authors approach \emph{higher helicities} via the Massey products under various assumptions about the vector fields or their domains (work in \cite{Gambaudo-Ghys97, Verjovsky94} concerns other approaches to the problem). In this paper we present an alternative to these approaches, which is a natural extension the notion of the linking number as a degree of a map, and exploits relations to the homotopy theory of certain maps associated to the link. The paper has two parts:

In the first part we show how the Milnor $\bar{\mu}_{123}$-invariant: $\bar{\mu}_{123}(L)$ of a parametrized borromean link $L$ in $S^3$ can be obtained as a Hopf degree of an associated map to the configuration space of three points in $S^3$. The presented approach to link-homotopy invariants of 3-component links has been proposed in \cite{Kohno02} for $n$-component links in $\R^3$, and conveniently simplified for 3-component links in $S^3$ in the joint work \cite{Deturck-Gluck-Melvin-Shonkwiler08}, where the full correspondence between $\bar{\mu}_{123}(L)$ and the Pontryagin-Hopf degree is proved. In Section \ref{sec:mu-and-hopf} we present the original proof of this correspondence in the borromean case, which is sufficient for our purposes.  

The second part of the paper is concerned with a new definition of the third order helicity denoted here by $\mathsf{H}_{123}(X;\mathcal{T})$, where $X$ is a volume preserving vector field having an invariant unlinked domain $\mathcal{T}\subset S^3$. The simplest of such domains are three invariant 
handlebodies in $S^3$ which have pairwise unlinked 
cycles in the first homology. This includes the case of borromean flux tubes already investigated in e.g. \cite{Berger90, Laurence-Stredulinsky00b}. In Section \ref{sec:ergodic}  we develop an ergodic formulation of $\mathsf{H}_{123}(X;\mathcal{T})$ as an average asymptotic $\bar{\mu}_{123}$-invariant in the spirit of Arnold's asymptotic linking number, which allows us to extend the definition of the invariant to topologically more complicated unlinked domains. We also indicate that $\mathsf{H}_{123}(X;\mathcal{T})$ provides a lower 
bound for the $L^2$-energy of $X$. 

%In Section \ref{sec:invariants-flows} we review Woltjer's helicity via the average asymptotic linking number, and use it as a springboard for defining a higher ``$\bar{\mu}_{123}$-helicity''. In Section \ref{sec:asymp-mu} we formulate an analog of higher helicity for a triple of vector fields and prove that it is an  invariant. This new invariant $\mathsf{H}_{123}(X,Y,Z)$ is defined on certain invariant domains of nonsingular volume preserving vector fields: $X$, $Y$, $Z$ in $S^3$, under the assumption that the field lines are asymptotically ``pairwise unlinked''. It may seem restrictive at first, but as it turns out, it is equivalent to the assumption on the pairwise helicities to vanish, and fits well into the the hierarchy of invariants proposed by Arnold and Khesin in \cite{Khesin98}. 

\bigskip
\no\emph{Acknowledgements:} My inspiration for the presented approach to the link-homotopy invariants  comes from the Toshitake Kohno paper, \cite{Kohno02}, whom I am grateful for the valuable e-mail correspondence. I have enjoyed conversations with many colleagues at the University of Pennsylvania, who have influenced this work. I wish to thank Herman Gluck for weekly meetings and his interest in this project, Frederic Cohen, Dennis DeTurck, Charlie Epstein, Paul Melvin, Tristan Rivi$\grave{\rm e}$re, Clay Shonkwiler, Jim Stasheff, David Shea Vick for the valuable input.  I am also grateful to my advisor 
Robert Ghrist who introduced me to the subject long time ago.

 The author acknowledges financial support from DARPA, \#FA9550-08-1-0386.



\section{The Milnor $\bar{\mu}_{123}$-invariant and the Hopf degree.}\label{sec:mu-and-hopf}

The $\bar{\mu}$-invariants of $n$-component links in $S^3$ have been introduced by Milnor in \cite{Milnor54, Milnor57}
as invariants of links up to \emph{link-homotopy}, recall that the link homotopy 
is a deformation of a link in $S^3$ which allows each component of the link to pass through 
itself but not through a different component. Clearly, this is a weaker equivalence than the equivalence of links up to \emph{isotopy} where components are not allowed to pass through themselves at all. The fundamental example of a $\bar{\mu}$-invariant is the linking number (denoted by $\bar{\mu}_{12}$) which is a complete invariant of the 2-component links up to link homotopy. In the realm of 3-component links the relevant invariants are the pairwise linking numbers $\bar{\mu}_{12}$, $\bar{\mu}_{23}$, $\bar{\mu}_{32}$, and the third invariant $\bar{\mu}_{123}$ in $\mathbb{Z}_{\text{gcd}(\bar{\mu}_{12}, \bar{\mu}_{23}, \bar{\mu}_{32})}$, which is 
a well defined integer if and only if $\bar{\mu}_{12}=\bar{\mu}_{23}=\bar{\mu}_{32}=0$. In the second part of the paper we will interpret this statement as a topological condition on the invariant set of a vector  field. A precise definition of $\bar{\mu}$-invariants is algebraic and involves the Magnus expansion of the lower central series of the fundamental group: $\pi_1(S^3-L)$ of the link complement. We refer the interested reader to the works in \cite{Milnor54, Milnor57}. In the remaining part of this section we will prove that $\bar{\mu}_{123}(L)$ is a Hopf degree for an appropriate map associated to the link $L$, provided that 
the link is \emph{borromean}, i.e. the pairwise linking numbers are zero (note that the borromean links are more general then so called Brunnian links, \cite{Milnor54}). 


Let us review basic facts about the Hopf degree $\mathscr{H}(f)$ of a map 
$f:S^3\longrightarrow S^2$, (see
e.g. \cite{Bott82}). The well known property of the Hopf degree is that  $\mathscr{H}: f\longrightarrow \mathscr{H}(f)$ provides an isomorphism between $\pi_3(S^2)$
and $\Z$. Recall that, up to a constant multiple, we
may express $\mathscr{H}(f)$ as ($M=S^3$)
\begin{gather}\label{eq:hopf-integral}
\mathscr{H}(f)= \int_{M} \alpha\wedge f^\ast\nu=\int_{M} \alpha\wedge \omega=\int_{M} \alpha\wedge d\alpha,
\end{gather}
where $\nu$ is the area 2-form on $S^2$, and $\alpha$ satisfies
$\omega=f^\ast\nu=d\alpha$. Notice that $f^\ast\nu$ is always exact
since the cohomology of $S^3$ in dimension $2$ vanishes (i.e.
$H^2(S^3)=0$). We may also interpret $\mathscr{H}(f)$ as an intersection number,
\cite{Bott82, Milnor97}. Namely, consider two regular values $p_1$ and $p_2\in S^2$ of the map $f$,
then $l_1=f^{-1}(p_1)$ and $l_2=f^{-1}(p_2)$ form a link in $S^3$, and the integral formula
\eqref{eq:hopf-integral} can be interpreted, as the intersection number of $l_1$ with the Seifert surface spanning $l_2$:
\begin{gather}\label{eq:hopf-intersection}
\mathscr{H}(f)=\text{lk}(l_1,l_2)\ .
\end{gather}
If we replace $S^3$ with an arbitrary closed
compact orientable 3-dimensional manifold $M$, we may still obtain
an invariant of $f:M\longrightarrow S^2$ this way, provided that
the condition $f^\ast\omega=d\alpha$ holds. 

\begin{proposition}\label{th:hopf-degree-general}
 Let $M$ be a closed Riemannian manifold, and $\nu\in \Omega^2(S^2)$ the area form on $S^2$.
 The formula \eqref{eq:hopf-integral} provides a homotopy invariant 
 for a map $f:M\longrightarrow S^2$, if the 2-form $f^\ast\nu$ is exact. Up to a constant multiple 
 this invariant can be calculated as an intersection number defined in \eqref{eq:hopf-intersection},
 where $l_1=f^{-1}(p_1)$ and $l_2=f^{-1}(p_2)$ form a link in $M$, where both $l_1$ and $l_2$ are null-homologous.
\end{proposition}

\begin{proof}
\no Given a homotopy $F:I\times M\mapsto
S^2$, $f_1=F(1,\,\cdot\,)$, $f_0=F(0,\,\cdot\,)$, we define
$\hat{\omega}=F^\ast\nu$. We have 
\begin{eqnarray*}
\hat{\omega} & = & F^\ast\nu=d\,\hat{\alpha},\\
\omega_1 & = & f_1^\ast\nu=i^\ast_1\,F^\ast\nu=d\,\alpha_1,\\
\omega_0 & = & f_0^\ast\nu=i^\ast_0\,F^\ast\nu=d\,\alpha_0,
\end{eqnarray*}
where $i_0:M\hookrightarrow M\times I$, $i_0(x)=(x,0)$, $i_1:M\hookrightarrow M\times I$, $i_1(x)=(x,1)$
are appropriate inclusions. The ``potentials'': $\hat{\alpha}\bigl|_{\{0\}\times M}$, $\alpha_0$, and $\hat{\alpha}\bigl|_{\{1\}\times M}$, $\alpha_1$ differ  by a closed form 
\[
\hat{\alpha}\bigl|_{\{0\}\times M}-\alpha_0=\beta_0,\quad \hat{\alpha}\bigl|_{\{1\}\times M}-\alpha_1=\beta_1,\quad d\beta_0=d\beta_1=0,
\]
therefore the Stokes Theorem immediately implies that Formula \eqref{eq:hopf-integral} is independent of the choice 
of the potential. For the proof of invariance we repeat the standard argument in \cite[p. 228]{Bott82}
\begin{eqnarray*}
 0 & = & \int_{F(M\times I)} \nu\wedge\nu
 =\int_{M\times I} \hat{\omega}\wedge \hat{\omega} =  \int_{M\times I} \hat{\omega}\wedge d\,\hat{\alpha}\\
 & = &\int_{M\times I} d(\hat{\omega}\wedge \hat{\alpha}) = \int_{\partial(M\times I)} \hat{\omega}\wedge \hat{\alpha}\ .
\end{eqnarray*}
Thus,
\begin{eqnarray*}
0 & = & \int_{\partial(M\times I)} \hat{\omega}\wedge
\hat{\alpha}  =  \int_{M} \omega_1\wedge \alpha_1- \int_{M}
\omega_0\wedge \alpha_0\\
& = & \int_{M} d\alpha_1 \wedge \alpha_1-\int_{M}
d\alpha_0\wedge \alpha_0 =  \mathscr{H}(f_1)-\mathscr{H}(f_0).
\end{eqnarray*}
\no The interpretation of $\mathscr{H}(f)$ as the intersection number \eqref{eq:hopf-intersection} is the same as in \cite[p. 230]{Bott82}.
\end{proof}


\no Given a 3-component parametrized link $L=\{L_1,L_2,L_3\}$ in $S^3$
we wish to associate a certain map $F_L:S^1\times S^1\times S^1\mapsto S^2$ to it, and interpret its Hopf degree as the Milnor $\bar{\mu}_{123}$-invariant. Recall, the definition of the configuration space of $k$ points in $M$:
\begin{gather*}
 \text{Conf}(M,k):=\{(x_1,x_2,\ldots, x_k)\in (M)^k\,|\,x_i\neq x_j,\text{for}\ i\neq j\}.
\end{gather*}
As an introduction to the method we review the Gauss formula for the linking number of 
a 2-component link $L=\{L_1, L_2\}$ in $\R^3$. Denote parameterizations 
of components by  $L_1=\{\mathsf{x}(s)\}$, $L_2=\{\mathsf{y}(t)\}$ and consider the map
\begin{gather*}
 F_L:S^1\times S^1\stackrel{L}{\longrightarrow} \Cnf(\R^3,2) \stackrel{r}{\longrightarrow} S^2,\qquad L(s,t)=(\mathsf{x}(s), \mathsf{y}(t))\ .
\end{gather*}
where $r(x,y)=\frac{x-y}{\|x-y\|}$ is a retraction of $\Cnf(\R^3,2)$ onto $S^2$.
It yields the classical Gauss linking number formula, up to a constant multiple:
\begin{gather*}
 \text{lk}(L_1,L_2)=\text{deg}(F_L),\qquad\quad \text{deg}(F_L)=\int_{S^1\times S^1} F^\ast_L(\nu),
\end{gather*}
where $\nu\in \Omega^2(S^2)$ is the 
area form on $S^2$. Consequently, the linking number $\text{lk}(L_1,L_2)$,
also known as the Milnor $\bar{\mu}_{12}$-invariant, can be obtained as the homotopy invariant 
of the map $F_L$ associated to $L$. Notice that homotopy classes $[S^1\times S^1,S^2]$ are 
isomorphic to $\Z$ and $\text{deg}:F\rightarrow \text{deg}(F)$ provides the isomorphism. In \cite{Kohno02},
Kohno proposed a natural extension of this approach to $n$-component parametrized links $L$ in $\R^3$ by considering maps $F_L:S^1\times\ldots\times S^1\longrightarrow \Cnf(\R^3,n)$ and their homotopy classes. More specifically, he proposed specific cohomology classes of the based loop space of $\Cnf(\R^3,n)$ as candidates for appropriate link-homotopy invariants of $L$. In \cite{Deturck-Gluck-Melvin-Shonkwiler08}, it has been observed that, in the $3$-component case it is  beneficial to consider $\Cnf(S^3,3)$, and $L\subset S^3$, since the topology of $\Cnf(S^3,3)$ simplifies dramatically (in comparison to $\Cnf(\R^3,3)$). We review this 
simplification in the following paragraph as it is essential for the proof of the next theorem.

Consider a 3-component link $L=\{L_1, L_2, L_3\}$ in $S^3$ 
parametrized by $\{\mathsf{x}(s),\mathsf{y}(t),\mathsf{z}(u)\}$ 
and the following map 
\begin{gather}\label{eq:F}
 F_L:S^1\times S^1\times S^1\stackrel{L}{\longrightarrow} \Cnf(S^3,3)\stackrel{H}{\longrightarrow} S^2,\qquad L(s,t,u)=(\mathsf{x}(s), \mathsf{y}(t), \mathsf{z}(u))\ .
\end{gather}
\no In order to explain the map $H:\Cnf(S^3,3)\longrightarrow S^2$, one must notice that 
$\Cnf(S^3,3)\subset S^3\times S^3\times S^3$ is
diffeomorphic to $S^3\times \Cnf(\R^3,2)=S^3\times
\R^3\times (\R^3\setminus\{0\})$, and consequently deformation
retracts onto $S^3\times S^2$. Denote by 
$H:S^3\times \R^3\times (\R^3\setminus\{0\})\mapsto S^2$ the projection on $S^2$ factor. 
Considering $S^3$ as unit quaternions, the map $H$ can be expressed explicitly by the formula, \cite{Deturck-Gluck-Melvin-Shonkwiler08}:
\begin{gather}\label{eq:hermans-map}
 \Cnf(S^3,3) \ni (x,y,z) \stackrel{H}{\longrightarrow}
 \frac{\text{pr}(x^{-1}\cdot y)-\text{pr}(x^{-1}\cdot z)}{\|\text{pr}(x^{-1}\cdot y)-\text{pr}(x^{-1}\cdot z)\|}\in S^2,
\end{gather}
where $\cdot$ stands for the quaternionic multiplication, $\ ^{-1}$
is the quaternionic inverse, and $\text{pr}:S^3\longrightarrow \R^3$ the
stereographic projection from $1$. As a result one has the following particular expression for $F_L$:
\begin{gather}\label{eq:F_L}
\notag F_L(s,t,u)=\frac{\text{pr}(\mathsf{x}(s)^{-1}\cdot \mathsf{z}(u))-\text{pr}(\mathsf{x}(s)^{-1}\cdot \mathsf{y}(t))}{\|\text{pr}(\mathsf{x}(s)^{-1}\cdot \mathsf{z}(u))-\text{pr}(\mathsf{x}(s)^{-1}\cdot
 \mathsf{y}(t))\|}.
\end{gather}

\no At this point we note that one has a freedom in choosing the deformation retraction $H$ in \eqref{eq:hermans-map} but this particular formula makes the proof of the main theorem of this section possible. By the \emph{borromean rings} we understand any $3$-component link with $\bar{\mu}_{123}=\pm 1$ link homotopic to a link with diagram presented on Figure \ref{fig:n-borromean} (where the sign can be determined from the orientation of components).

\begin{theorem}\label{th:milnor-hopf}
 Let $L^\text{Bor}=\{L_1,L_2,L_3\}$ be a 3-component link in $S^3$ link homotopic to the Borromean rings, and consider
 the associated map $F_{L^\text{Bor}}:S^1\times S^1\times S^1\longrightarrow S^2$ defined in \eqref{eq:F}. Then the Hopf degree of the map 
 $F_{L^\text{Bor}}$ satisfies
 \begin{gather}\label{eq:hopf=2}
   \mathscr{H}(F_{L^\text{Bor}})=\pm 2,
 \end{gather}
 where the sign depends on the choice of orientations.
\end{theorem}

\begin{proof}

\no According to Proposition \ref{th:hopf-degree-general}, $\mathscr{H}(F_L)$ is well defined for a link $L\subset S^3$ provided  $F_L^\ast\nu\in \Omega^2(\mathbb{T})$ is trivial in $H^2(\mathbb{T})$, $\mathbb{T}=S^1\times S^1\times S^1$. This fact will follow from the proof, but we will show in Corollary \ref{th:pairwise-linking}, that it holds if and only if the pairwise linking numbers of $L$ are zero. The method relies on a direct calculation of 
$\mathscr{H}(F_{L^\text{Bor}})$, for a carefully chosen parametrization of $L^\text{Bor}$ in $S^3$. This calculation is achieved by 
visualization of the link $l_{S,N}=l_S\cup l_N$ in $\mathbb{T}$, and application of Formula \eqref{eq:hopf-intersection}, where
\[
 l_S:=F^{-1}_{L^\text{Bor}}(S),\qquad l_N:=F^{-1}_{L^\text{Bor}}(N)
\]
are preimages of the North pole $N=(0,0,1)$ and South pole $S=(0,0,-1)$ in $S^2$. Notice that $[F_L^\ast\nu]=0$ in $H^2(\mathbb{T})$ if and only if $[l_S]=0$ and $[l_N]=0$ in $H_1(\mathbb{T})$.


 We identify $S^3$ with the set of unit quaternions in $\R^4$ with standard
coordinates: 
\[
 (w,x,y,z)=w+x\,\mathbf{i}+y\,\mathbf{j}+z\,\mathbf{k},
\]
and choose a specific parametrization of the
Borromean rings $L^\text{Bor}$ in $S^3$, as follows: define the $L_1$ component of
 $L^\text{Bor}$ to be the great circle through $1$ in $S^3$, parametrized as
\[
\mathsf{x}(s)=\cos(s)+\sin(s)\,\mathbf{k},\quad
\bigl(\,\mathsf{x}(s)^{-1}=\cos(s)-\sin(s)\,\mathbf{k}\,\bigr).
\]
Notice that $\text{pr}(\mathsf{x}(s))$ parameterizes the $z$-axis in $\R^3$. Because
\begin{eqnarray*}
\mathsf{x}(s)^{-1}\cdot(w,x,y,z) & = & \bigl(\cos(s) w+\sin(s) z, \cos(s) x +\sin(s) y,\\
& & \ \cos(s) y-\sin(s) x, \cos(s) z-\sin(s) w\bigr),
\end{eqnarray*}
we observe that multiplication by $x(s)^{-1}$ is a rotation by angle $s$
in $(w,z)$-plane and $(x,y)$-plane of $\R^4$ , and the flow defined
by this $S^1$-action is tangent to the great circles of $S^3$. 
Consequently, the projected flow on $\R^3$
via the stereographic projection $\text{pr}:S^3\subset \R^4\longrightarrow \R^3$
presents the standard picture of the Hopf fibration, and the $S^1$-action ``pushes'' points in $\R^3$ along Hopf circles. 

\begin{figure}
\begin{center}
\begin{picture}(414,160)
  \put(15,0){\includegraphics{borromean3}}
  \put(80,120){$L_1$}
  \put(35,70){$L_2$}
  %\put(360,60){$L_3$}
  \put(60,5){$L_3$}
  \put(315,80){$A_2$}
  \put(240,60){$A_1$}
\end{picture}  
\caption{The model of $L_\text{Borr}$ parametrized by $\{L_1=\text{pr}(\mathsf{x}(s)),L_2=\text{pr}(\mathsf{y}(t)),L_3=\text{pr}(\mathsf{z}(u))\}$. Arc $A_1$ is a part of the unit circle on $xy$-plane, $A_2$ is a part of circle of radius $1+\epsilon$.} \label{fig:conf1}

\end{center}
\end{figure}

Let us call a Hopf torus in $\R^3$ the $r$-torus if and only if it
contains a circle of radius $r$ in the $(x,y)$-plane and is
invariant under the $S^1$-action.
 Figure \ref{fig:conf1} shows how to define the second and the third component $\{L_2, L_3\}$ of the Borromean rings $L_\text{Borr}$ in the stereographic projection.
$L_2$ will be the annuli with a rounded wedge removed, i.e. an arc $A_1$ of the circle
of radius $1$ and $A_2$ arc belongs to the circle of radius $r_\epsilon=(1+\epsilon)$ in the $(x,y)$-plane. 
%Notice that $A_1$ is a part of the Hopf circle
%(i.e. the unit circle) and therefore invariant under the $S^1$-action. The arc $A_2$ belongs to the
%$\frac{1}{2} r_\epsilon$-torus.
 The component $L_3$ is chosen to be a vertical ellipse
linking with $L_2$, which belongs to the $r_{\epsilon/2}$-torus, and
projects on the $(x,y)$-plane as a segment. Recall that every point on an $r$-torus traces a
$(1,1)$-curve under the $S^1$-action. This motion, for sufficiently small $\epsilon$, can be regarded 
as a composition
of the rotation by angle $s$ in both the direction of the meridian
and the longitude of an $r$-torus.  Therefore for different values of $s$ the $S^1$-action ``rotates'' the components $L_2$ and $L_3$, by sliding 
along the Hopf tori by angle $s$ in the meridian and the longitudinal
direction, we denote links obtained this way by 
\[
 L^s_2(t)=\text{pr}(\mathsf{x}^{-1}(s)\cdot\mathsf{y}(t)),\qquad \text{and}\qquad  L^s_3(u)=\text{pr}(\mathsf{x}^{-1}(s)\cdot\mathsf{z}(u)).
\]
Let us collect useful information below
\begin{itemize}
\item[$(i)$] The $z$-axis stays invariant under the $S^1$-action.
\item[$(ii)$] $A_2$ slides, under the $S^1$-action, on the $r_\epsilon$-torus, $A_1$ simply
rotates around the $z$-axis in the $(x,y)$-plane.
\item[$(iii)$] We can disregard the longitudinal rotation (i.e.
pretend that the observer sits on the unit circle) the motion of arcs:
$L^s_3$, $A^s_2$ together with the connecting segments is simply a
full rotation around the unit circle on the $(x,y)$-plane (which can be thought of as the ``axis of rotation'').
Stages of this motion are pictured on the lower part of Figure \ref{fig:4-pos}.
\end{itemize}



  We seek to visualize the projection of $l_{S,N}=l_S\cup l_N$ on the $su$-face and $st$-face of the domain $\mathbb{T}=S^1\times
S^1\times S^1$ of $F_{L^\text{Bor}}$ parameterized by $(s,t,u)\in \mathbb{T}$, (it is convenient to  think about $\mathbb{T}$ as a cube in $(s,t,u)$-coordinates, see Figure \ref{fig:framed-link}). 

For example when $s=0$, ($\mathsf{x}(0)=1$), a point
$(0,t_0,u_0)$ belongs to $l_N$, if and only if the vector $\mathsf{v}_0=\text{pr}(\mathsf{y}(t_0))-\text{pr}(\mathsf{z}(u_0))$
``points North'', i.e. in the direction of $N=(0,0,1)$, analogous condition holds for 
$S=(0,0,-1)$, and $l_S$. Thus in order to determine a diagram of $l_{S,N}$, we must keep track of the ``head'' and ``tail'' of the vector $\mathsf{v}_s=L^s_2(t)-L^s_3(u)$, for various values of $s$ and record values of $t$ and $u$
for which $\mathsf{v}_s$ points ``North'' and ``South''
(Figure \ref{fig:4-pos}). We have the following condition
\begin{equation}\label{eq:arrow-eq}
  (s,t,u)\in l_{S,N}\qquad\text{if and only if},\qquad \mathsf{v}_s \parallel S\ \text{or}\ N.
\end{equation}
 An alternative way to view this condition is the following: for a given value of $s$ 
the coordinates $u$ and $t$ of points in $l_{S,N}$ can be identified with points on $L^s_2$
and $L^s_3$, which come from the intersection points in $\pi(L^s_3\cup L^s_2)$, under the projection $\pi(x,y,z)=(x,y)$, onto the $(x,y)$-plane (see the bottom part of Figure \ref{fig:4-pos}). These intersection points are simply crossings in the diagram of the link 
$\{L^s_2,L^s_3\}$ on the $(x,y)$-plane. Points $N$ and $S$ can be assumed to be regular values of $F_{L_\text{Borr}}$ because $\pi(L^s_2(t))$ and $\pi(L^s_3(u))$ always intersect transversely (bottom of Figure \ref{fig:4-pos}), which is persistent under small perturbations of the parametrization of $L_\text{Borr}$.

The main task ahead is to visualize $l_N$ and $l_S$. For a fixed $s=s_0$ we draw arrows:
$N$ and $S$ (Figure \ref{fig:4-pos}) which lie directly above the crossings of the diagram of 
$\{L^{s_0}_2,L^{s_0}_3\}$ on the $(x,y)$-plane. Next, we trace each component $L^{s_0}_2$ 
and $L^{s_0}_3$ and record values of $t$ and $u$ parameter for which the condition \eqref{eq:arrow-eq} holds, without loss of generality we assume that $L^{s}_2$ is parametrized by the unit $t$-interval, and $L^{s}_3$ is parametrized by the unit $u$-interval.
\begin{figure}[htbp]
\begin{center}
\begin{picture}(498, 590)
  \put(30,20){\includegraphics{4-pos}}
  \put(30,260){$s=\frac{3\pi}{2}$}
  \put(30,360){$s=\pi$}
  \put(30,460){$s=\frac{\pi}{2}$}
  \put(30,550){$s=0$}
\end{picture}
\caption{Four positions corresponding to angles: $s=0,\frac{\pi}{2}, \pi, \frac{3\pi}{2}$, small arrows next to $N$ and $S$ indicate a motion as $s \nearrow$. Several phases of the rotation by $s$ are pictured on the bottom four drawings.} \label{fig:4-pos}
\end{center}
\end{figure}
 This process is elucidated below for $s=0, \frac{\pi}{2}, \pi,
\frac{3\pi}{2}$, which is sufficient to draw projections of $l_S$ and $l_N$ (each position is pictured on Figure \ref{fig:conf1}) on $st$ and $tu$-faces of $\mathbb{T}$. We
describe each case separately,
\begin{itemize}
\item[($s=0$):] This is the initial position, where $L^0_3$ is
vertical. We mark the initial point $L^0_3(0)$ on the top of $L^0_3$
and orient $L^0_3$ counter-clockwise, the initial point $L^0_2(0)$ on
$L^0_2$ is chosen in the middle of the left segment connecting
$A_1$ and $A_2$, and the middle point $L^0_2(\frac{1}{2})$ is located
the middle of the left segment connecting $A_1$ and $A_2$. Points in
$l_S$  are marked by $\times$ and points in $l_N$,
 are marked by $\bullet$. In this configuration there are
four vertical arrows, two pointing North and two South, this gives
us four points along the $t$-interval ordered as $NSSN$, ($|-\bullet-\times-\times-\bullet-|$) and four points along the $u$-interval
ordered ($|-(\times\bullet)--(\bullet\times)-|$), where $(\times\bullet)$ 
indicates the ``collision'' i.e. common value of $u$ for both $l_S$ and $l_N$.
\end{itemize}
\no The remaining cases are self-explanatory, see Figure
\ref{fig:conf1},
\begin{itemize}
\item[($s=\frac{\pi}{2}$):] $L^s_3$ is horizontal,
and $t$-coordinates of $l_{S,N}$ are colliding
($|-(\bullet\times)-(\times\bullet)-|$). $L^s_2$ is vertical and
$u$-coordinates of $l_{S,N}$ are ordered as  $(NNSS)$,
($|-\bullet-\times-\times-\bullet-|$).
\item[($s=\pi$):] $L^s_3$ is vertical again, and $t$-coordinates of $l_{S,N}$ are ordered
: $SNNS$, ($|-\times-\bullet-\bullet-\times-|$). $L^s_2$ is
horizontal and $u$-coordinates of $l_{S,N}$ are colliding,
($|-(\bullet\times)--(\times\bullet)-|$).
\item[($s=\frac{3\pi}{2}$):] $L^s_3$ is horizontal,
and $t$-coordinates of $l_{S,N}$ are colliding
($|-(\times\bullet)-(\bullet\times)-|$). $L^s_2$ is vertical and
$u$-coordinates of $l_{S,N}$ are ordered as $(NNSS)$,
($|-\times-\times-\bullet-\bullet-|$).
\end{itemize}

\begin{figure}[htbp]
\begin{center}
\begin{picture}(480, 580)
  \put(140,420){\includegraphics[width=3in]{cube2-dashed}}
  \put(40,20){\includegraphics[width=5in]{proj-red-green-dashed}}
  \put(30,180){$s$}
  \put(30,390){$s$}
  \put(227,180){$s$}
  \put(227,390){$s$}
  \put(30,270){$\frac{\pi}{2}$}
  \put(227,270){$\frac{\pi}{2}$}
  \put(200,10){$u$}
  \put(313,8){(D)}
  \put(115,8){(C)}
  \put(400,10){$t$}
  \put(200,220){$u$}
  \put(115,215){(A)}
  \put(400,220){$t$}
  \put(313,215){(B)}
\end{picture}
\caption{Projection of $l_{S,N}$ on the $su$-face and $st$-face of $\mathbb{T}$. The strands $l_S$ and $l_N$ are oppositely oriented since $l_{S,N}$ is null homologous in $\mathbb{T}$.}
\label{fig:framed-link}
\end{center}
\end{figure}

\no Collecting this information, we may draw the projection of
the link $l_{S,N}$ on the $su$-face of $\mathbb{T}$, which is represented
by square (A) in Figure \ref{fig:framed-link}. Analogously, the projection of 
$l_{S,N}$ on the $st$-face of $\mathbb{T}$,  is represented
by square (B).
The dashed line indicates the $l_N$ component of 
$l_{S,N}$ and the solid line indicates $l_S$. 
In order to obtain the diagram of
$l_{S,N}$ we will resolve the double points on Diagram (A) into crossings, 
((B) can be resolved analogously). For example let us resolve the 
``circled'' double point on (A), which occurs for $s=\frac{\pi}{2}$, of the \emph{left} two stands of $l_{S,N}$. It suffices to determine the value of the $t$-coordinate at this point.
Diagram (B) tells us that $l_S$ is below $l_N$ because $\mathbb{T}$ is oriented so that the $t$-axis points above the $su$-face (see top of Figure \ref{fig:framed-link}).
 Resolving the remaining crossings in
a similar fashion leads to a diagrams of $l_{S,N}$ presented in squares (C) and (D). 
Diagram (C) is easier to use for the calculation of $\mathscr{H}(F_{L_{\text{Bor}}})=\text{lk}(l_S, l_N)$ as an intersection number of 
e.g. $l_S$ with a Seifert surface spanning $l_N$.
 Observe that strands in $l_S$
and $l_N$ are oppositely oriented, which follows from the fact that
they are null-homologous in $H^1(\mathbb{T})$, there are an obvious annuli bounding $l_S$ and $l_N$ 
on Diagram (C). Clearly, the linking number of
$l_S$ and $l_N$ is equal to $\pm 2$, since the intersection number of e.g. $l_S$ with the obvious annulus 
on Diagram (C) is $\pm 2$. This proves \eqref{eq:hopf=2}.
\end{proof}

\begin{figure}[htbp]
\centering
\begin{picture}(380, 150)
  \put(0,0){\includegraphics[width=5.2in]{borr}}
  \put(280,7){$\mathbf{\Bigl\}}$}
  \put(285,22){$n$}
\end{picture}  
\caption{left: $\bar{\mu}_{123}=\pm 1$ and right: $\bar{\mu}_{123}=\pm n$.}
\label{fig:n-borromean}
\end{figure}



\begin{corollary}\label{th:general-mu_123}
 If $L=\{L_1,L_2,L_3\}$ is a 3-component link with vanishing pairwise linking numbers then
 \begin{gather*}
  \mathscr{H}(F_L)=\pm 2\,\bar{\mu}_{123}
 \end{gather*}
\end{corollary}
\begin{proof}
 Clearly, if $L$ and $L'$ are link-homotopic then $F_{L}$ and $F_{L'}$ are homotopic maps.
 By the Milnor classification of 3-component links up to link homotopy (see \cite{Milnor54}), every 3-component link $L$ with zero pairwise linking numbers and $\bar{\mu}_{123}=\pm n$ is represented by the right diagram on Figure \ref{fig:n-borromean}. Consequently, up to homotopy, the associated map $F_L$ can be obtained from $F_{L_\text{Borr}}$ by covering one of the $S^1$ factors in $\mathbb{T}$, $n$-times. 
\end{proof}
\begin{lemma}[\cite{Cohen-private}]
 We have the following isomorphisms
\begin{gather*}
\pi_2 (\Cnf(\R^3,3))\cong \Z\oplus \Z\oplus \Z,\qquad \pi_2
(\Cnf(S^3,3))\cong \Z,
\end{gather*}
and the induced homomorphism $i^\ast:\pi_2 (\Cnf(\R^3,3))\mapsto
\pi_2(\Cnf(S^3,3))$, from $i:\R^3\mapsto \R^3\cup \{\infty\}\cong S^2$, maps the generators $A_{ij}$, $1\leq i\leq
j\leq 3$, to $\pm\,\alpha$, where $\alpha$ is the single generator
of $\pi_2(\Cnf(S^3,3))$.
\end{lemma}
\begin{proof}
The Hurewicz Theorem tells us that $H_2(\Cnf(\R^3,3))\cong
\pi_2(\Cnf(\R^3,3))$. The 2-skeleton of $\Cnf(\R^3,3)$ has a homotopy
type of $S^2\vee S^2\vee S^2$ and $H_2(\Cnf(\R^3,3))$ is generated by 
$A_{ij}: S^2\mapsto \Cnf(\R^3,3)$, $1\leq i<j\leq 3$ representing each 2-sphere in the 2-skeleton (c.f. \cite{Cohen02}).
The group of permutations $\mathbb{S}_3$ acts on $\Cnf(\R^3,3)$ by permuting coordinates.
As a consequence of this we obtain the action on generators of $H_2(\Cnf(\R^3,3))$ by permutation.
We have seen that $\Cnf(S^3,3)$ has the homotopy type of $S^3\times S^2$ thus $\pi_2(\Cnf(S^3,3))=\Z$, and let us  denote the generator by $\alpha$. 
Now, $i^\ast{A_{ij}}=k\,\alpha$, $k\in \Z$, notice that $i^\ast{A_{23}}=\alpha$ since 
 $i\circ A_{23}=\text{id}_{S^2}$, using the action by permutations we obtain 
 $ i^\ast{A_{ij}}=\alpha $.
\end{proof}

\begin{corollary}\label{th:pairwise-linking}
Denote by $F^{kl}_L=F_L\bigl|_{S^1_k\times S^1_l}$ restriction of $F_L$ defined in \eqref{eq:F}
to the $kl$-torus $k,l\in\{s,t,u\}$, $k\neq l$. Then
\begin{gather}
 \text{deg}(F^{kl}_L)=\text{lk}(L_k,L_l),
\end{gather}
where $\text{lk}(L_k,L_l)$ is the linking number between the components
parameterized by $k$, and $l$.
\end{corollary}
\begin{proof}
 By definition 
 \[
  \text{deg}(S_i^1\times S_j^1\stackrel{L_{ij}}{\longrightarrow} \Cnf(\R^3,3)\stackrel{\pi_{ij}}{\longrightarrow} \Cnf(\R^3,2)\cong S^2)=\text{lk}(L_i,L_j),
  \] 
thus on the cohomology level $\text{lk}(L_i,L_j)=L^\ast_{ij}[A_{ij}]=(F^{ij}_L)^\ast(\alpha)$.
\end{proof}
 
\no Results of Theorem \ref{th:milnor-hopf} and Proposition \ref{th:hopf-degree-general} combined with Formula \eqref{eq:hopf-integral} allow us to express $\bar{\mu}_{123}(L)$ of a borromean link $L$ as
 \begin{eqnarray}\label{eq:mu-formula}
  \bar{\mu}_{123}(L) & = & \int_{\mathbb{T}} F^\ast_L\nu\wedge \alpha=\int_{\mathbb{T}} L^\ast\omega\wedge \alpha,\\
  \notag & \text{for} &  \quad F^\ast_L\nu=d\alpha,\quad \omega=H^\ast\nu\in \Omega^2(\Cnf(S^3,3)),
 \end{eqnarray}
 where $\nu$ is the area form on $S^2$, and $H:\Cnf(S^3,3)\longrightarrow S^2$ is the deformation retraction 
 (as e.g. in \eqref{eq:hermans-map}). Alternatively, we may view $\omega$ as a 2-form on $(S^3)^3$ which is singular along 
 the diagonals $\mathbf{\Delta}\subset (S^3)^3$, and the singularity is of order $O(r^2)$, where $r$ is a distance to $\mathbf{\Delta}$. Consequently, $\omega$ is integrable but not square integrable on $(S^3)^3$. 
 We also point out that the 2-form $\omega$ can be replaced with any other cohomologus 2-form in $\text{Conf}(S^3,3)$.
 
\begin{remark}\label{rem:anti-symmetry}
  Notice that the integral formula \eqref{eq:mu-formula} is antisymmetric in the following sense 
  \begin{eqnarray*}
   \bar{\mu}_{123}(L_1,L_2,L_3) & = & \int_{S^1_s}\int_{S^1_t}\int_{S^1_u} \omega\wedge\alpha(\dot{\mathsf{x}}(s),\dot{\mathsf{y}}(t),\dot{\mathsf{z}}(u))\,ds\wedge
   dt\wedge du\\
   & = & \int_{S^1_s}\int_{S^1_t}\int_{S^1_u} -\omega\wedge(-\alpha)(\dot{\mathsf{y}}(t),\dot{\mathsf{x}}(s),\dot{\mathsf{z}}(u))\,ds\wedge dt\wedge du\\
   & = & \int_{S^1_t}\int_{S^1_s}\int_{S^1_u} \omega\wedge\alpha(\dot{\mathsf{y}}(t),\dot{\mathsf{x}}(s),\dot{\mathsf{z}}(u))(-\,dt\wedge ds \wedge du)=-\bar{\mu}_{123}(L_2,L_1,L_3)
  \end{eqnarray*}
  where in the second identity the potential changes sign because the domain changes orientation.
\end{remark}

\no \section{Invariants of volume preserving flows. Helicity.}\label{sec:invariants-flows}

Given finitely many  volume preserving vector fields $X_1$, $X_2$, $\ldots X_k\in \text{Vect}^\text{vol}(M)$  on $M=S^3$ or a homology 3-sphere one seeks quantities $\mathsf{I}(X_1,X_2,\ldots,X_k)$ invariant under the action of volumorphisms isotopic to the identity $g\in \text{Diff}^\text{vol}_0(M)$, commonly known as \emph{higher-helicities}:
\begin{gather}\label{eq:invariants}
 \mathsf{I}(X_1,X_2,\ldots,X_k)=\mathsf{I}(g_\ast X_1,g_\ast X_2,\ldots,g_\ast X_k),\qquad\quad \text{for all }g\in \text{Diff}^\text{vol}_0(M), 
\end{gather}
where $g_\ast$ is a push-forward by a diffeomorphism $g$. To distinguish the case of a single vector field $X$ (i.e. $X=X_1=\ldots=X_k$) we often refer to $\mathsf{I}(X)=\mathsf{I}(X_1,X_2,\ldots,X_k)$ as \emph{self-helicity}. We elucidated in the introduction a fundamental example of such invariant is the ordinary helicity $\mathsf{H}(X_1,X_2)$ of a pair of vector fields. In the the remaining part of this section we review 
well known formulations of the helicity, which will later help us to point out analogies to the 
proposed formulation of the 3rd order helicity. Let $\mathcal{T}=\mathcal{T}_1\cup \mathcal{T}_2$ represent two invariant subdomains (not necessarily disjoint) under flows of $X_1$ and $X_2$ in $S^3$ and let 
\[
 \mathcal{T}=\mathcal{T}_1\times\mathcal{T}_2\subset \Cnf(M,2)\subset M\times M.
\]
Recall that the formula for $\mathsf{H}(X_1,X_2)$, from \cite{Khesin05, Vogel03}, specialized to invariant subdomains $\mathcal{T}=\mathcal{T}_1\cup \mathcal{T}_2$ may be expressed as 
\begin{equation}\label{eq:helicity}
 \mathsf{H}_{12}(X_1,X_2)=\int_{\mathcal{T}_1\times\mathcal{T}_2} \omega\wedge\iota_{X_1}\mu\wedge\iota_{X_2}\mu,
\end{equation}
where $\omega$ is so called linking form on $M\times M$. When $\mathcal{T}=M\times M$ this formula is equivalent to a more commonly known expression: $\mathsf{H}(X_1,X_2)=\int_M \iota_{X_1}\mu\wedge d^{-1}(\iota_{X_2}\mu)$ (because $\omega$ also represents the integral kernel of $d^{-1}$).

 Philosophically, $\mathsf{H}(X_1,X_2)$  can be derived (c.f. \cite{Arnold86}) from the linking number of a pair of closed curves, which is expressed by Arnold's Helicity Theorem. For orbits
 $\{\mathscr{O}^1(x),\mathscr{O}^2(y)\}$  of $X_1$ and $X_2$ through $x, y\in M$, we introduce the following notation for the long pieces of closed up orbits
 
\begin{eqnarray}\label{eq:orbits}
 \mathscr{O}^i_T(x) & = & \{\Phi^i_t(x)\ |\ 0\leq t\leq T\}\subset \mathscr{O}^i(x),\qquad i=X, Y,\\
 \bar{\mathscr{O}}^i_T(x) & := & \mathscr{O}^i_T(x)\cup \sigma(x,\Phi^i(x,T)).
\end{eqnarray}
 where $\sigma(x,y)$ denotes a short path, \cite{Vogel03}, connecting $x$ and $y$ in $M$ (see Section \ref{sec:ergodic}).
 
\begin{theorem}[Arnold's Helicity Theorem, \cite{Arnold86}]
 Given $X_1, X_2\in\text{Vect}^\text{vol}(M)$,  The following 
 limit exists almost everywhere on $M\times M$:
 \begin{gather}\label{eq:linking-function}
  \bar{m}_{X_1 X_2}(x,y)=\lim_{T\to \infty} \frac{1}{T^2}\, \text{lk}(\bar{\mathscr{O}}^1_T(x)\times\bar{\mathscr{O}}^2_T(y))\ .
 \end{gather}
Moreover, $\bar{m}_{X_1 X_2}$ is in $L^1(M\times M)$, and 
\begin{equation}\label{eq:helicity-asymp}
 \mathsf{H}_{12}(X_1,X_2)=\int_{M\times M} \bar{m}_{X_1 X_2}(x,y)\, \mu(x)\wedge\mu(y)\ .
\end{equation}
\end{theorem}
\no The function $\bar{m}_{X_1 X_2}$ represents an asymptotic linking number of orbits $\{\mathscr{O}^1(x),\mathscr{O}^2(y)\}$, and the identity \eqref{eq:helicity-asymp} tells us that 
the helicity $\mathsf{H}_{12}(X_1,X_2)$ is equal to the average asymptotic linking number. We will demonstrate, in coming paragraphs, how this philosophy is applied to obtain the asymptotic $\bar{\mu}_{123}$-invariant for 
3-component links and the third order helicity.

\section{Definition of ``$\bar{\mu}_{123}$-helicity'' on invariant unlinked handlebodies.}\label{sec:handlebody}

\no In this section we apply the formulation of the $\bar{\mu}_{123}$-invariant for the 3-component links in $S^3$, 
obtained in Section \ref{sec:mu-and-hopf}, to define the third order helicity of a volume preserving vector field $X$ on certain invariant domains $\mathcal{T}$ in $S^3$. 
As a ``warm-up'' to a more general case presented in Section \ref{sec:ergodic}, in the following paragraphs we consider the case of three disjoint unlinked handlebodies: $\mathcal{T}=\mathcal{T}_1\cup 	\mathcal{T}_2\cup \mathcal{T}_3$ each of genus $g(\mathcal{T}_i)$ in $S^3$.
When $\mathcal{T}$ represents three unlinked tubes (so called \emph{flux tubes}, \cite{Khesin98}) in $\R^3$, the third order helicity has been developed by several authors \cite{Berger90,Mayer03,Laurence-Stredulinsky00b} via the Massey product formula for the $\bar{\mu}_{123}$-invariant, we compare our approach to these known works in Section \ref{sec:massey}.
%This is slightly more general then the case of \emph{flux tubes},  modeled   a 
%link $L$, i.e. when $\mathcal{T}$ deformation retracts onto $L=L_1\cup L_2\cup L_3$, where $L_j$'s are the core curves of solid tori $\mathcal{T}_j$. 


 We consider $\mathcal{T}_i$ to have smooth boundary and the vector field $X$ to be tangent to  $\partial \mathcal{T}_i$, set
\[
 X_i:=X\bigl|_{\mathcal{T}_i},\qquad i=1,2,3,
\]
 denote the flow of $X$ on $S^3$ by $\Phi$, and flows of restrictions $X_i$ by $\Phi^i$.  Clearly,
 $\mathcal{T}$ is an invariant set of $X$.
 Given for any domain $\mathcal{T}$ with three connected components $\{\mathcal{T}_i\}$ we may always associate a product domain in $\Cnf(S^3,3)$ as follows
\begin{equation}\label{eq:admissible-domain}
 \mathcal{T}:=\mathcal{T}_1\times \mathcal{T}_2\times \mathcal{T}_3\subset \text{Conf}(S^3,3)\subset S^3\times S^3\times S^3\ .
\end{equation}
Notice that $\mathcal{T}$ is a domain with corners in $\Cnf(S^3,3)$, and we use the same notation for the product of $\mathcal{T}_i$ as for the union in $S^3$. Wherever needed we also assume that 
$(S^3)^3$ is equipped with a product Riemannian metric.
 A domain $\mathcal{T}$ defined in \eqref{eq:admissible-domain}, where $\mathcal{T}_i\cap \mathcal{T}_j=\O$, $i\neq j$ and each $\mathcal{T}_i$ is a handlebody in $S^3$ is called \emph{unlinked handlebody} if and only if the 2-form $\omega\in\Omega^2(\Cnf(S^3,3))$ defined in Equation \eqref{eq:mu-formula} is exact on $\mathcal{T}$, i.e. $\omega$ has a local \emph{potential} $\alpha_\omega\in \Omega^1(\mathcal{T})$: 
\begin{equation}\label{eq:omega-exact}
 \omega=d\alpha_\omega\ .
\end{equation}
Denote a volume preserving vector field and an unlinked handlebody $\mathcal{T}$ as $(X;\mathcal{T})$.


\begin{remark}
\no Notice that since $\omega$ is a dual cohomology class to the $S^2$ factor of $\Cnf(S^3,3)\cong S^3\times S^2$, thus $\omega$ does not admit a global potential.
\end{remark}
\no Because each handlebody $\mathcal{T}_i$ has a homotopy type of a bouquet of circles $H_1(\mathcal{T}_i)$ is generated by cycles $\{L^k_i\}_{k=1,\ldots,g(\partial\mathcal{T}_i)}$ corresponding to the circles, which we call \emph{basic cycles}. We have the following practical characterization of unlinked handlebodies:

\begin{lemma}
  $\mathcal{T}$ is an unlinked handlebody if and only if 
  \begin{equation}\label{eq:admissible-condition}
  \text{lk}(L^k_i,L^r_j)  = 0,\qquad \forall\, i,j,\qquad i\neq j
 \end{equation}
\end{lemma}
\begin{proof}
The standard pairing, \cite{Bott82}, $H^2(\mathcal{T})\times H_2(\mathcal{T})\longrightarrow \R$ implies that a closed $k$-form is exact if and only if it evaluates to zero on all $k$-cycles of the domain. By the K$\ddot{\rm u}$nneth formula $H_2(\mathcal{T})=H_2(\mathcal{T}_1\times \mathcal{T}_2\times \mathcal{T}_3)$ is generated by $L^k_i\otimes L^r_j$. Corollary \ref{th:pairwise-linking} implies (for $i\neq j$) that
\[
 \text{lk}(L^k_i,L^r_j) = \int_{L^k_i\times L^r_j} \omega\ .
\]
Therefore the condition \eqref{eq:admissible-condition} is necessary and sufficient for $\omega$ to be exact on $\mathcal{T}$.
\end{proof}

\begin{figure}[htbp]
	\centering
		\includegraphics[width=\textwidth]{borromean-tube2}
	\caption{The simplest unlinked handlebody: flux tubes $\mathcal{T}^\text{Borr}$ modeled on the borromean rings.}
	%\caption{Unlinked handlebodies.}
	\label{fig:invariant-sets}
\end{figure}

\no Let $X$ be a volume preserving vector field on $S^3$ which admits an unlinked domain $\mathcal{T}$ which 
is invariant under its flow. We define the $\bar{\mu}_{123}$-helicity of $(X;\mathcal{T})$  denoted by $\mathsf{H}_{123}(X;\mathcal{T})$ or $\mathsf{H}_{123}(X_1,X_2,X_3)$
as follows:
\begin{equation}\label{eq:mu_123-helicity}
\boxed{\mathsf{H}_{123}(X;\mathcal{T})=\mathsf{H}_{123}(X_1,X_2,X_3)\stackrel{\text{def.}}{=}\int_{\mathcal{T}} \bigl(\alpha_\omega\wedge d\alpha_\omega\bigr)\wedge \iota_{X_1}\mu_1\wedge \iota_{X_2}\mu_2\wedge \iota_{X_3}\mu_3, }
\end{equation}
where $\mu_i$ denotes the pull-back of the volume form $\mu$ on $S^3$ under the projection 
\begin{equation}\label{eq:projection-i}
\pi_i:S^3\times S^3\times S^3\longrightarrow S^3,\qquad \pi_i(x_1,x_2,x_3)=x_i,
\end{equation}
and $\iota_{X_i}$ is a contraction by a vector field $X_i$, 
 our notational convention is to denote by $\iota_{X_i}\mu_i$ both the forms on the base of $\pi_i$ and the pullbacks: $\pi^\ast_i\bigl(\iota_{X_i}\mu_i\bigr)$. Notice that $\mu=\mu_1\wedge\mu_2\wedge\mu_3$ is a natural volume form on the product: $S^3\times S^3\times S^3$.  
 
 There are obvious analogies between the formula above, the formula for the helicity $\mathsf{H}_{12}(X;\mathcal{T})$ and the integral formula for the 
$\bar{\mu}_{123}$-invariant \eqref{eq:mu-formula}. Clearly, the 3-form 
 \[
 \gamma_\omega:=\alpha_\omega\wedge d\alpha_\omega=\alpha_\omega\wedge \omega,
 \]
 plays a role of the linking form as $\omega$ in the formula \eqref{eq:helicity} for $\mathsf{H}_{12}(X;\mathcal{T})$. 
 The main motivation behind this definition is the ergodic interpretation  of $\mathsf{H}_{123}(X;\mathcal{T})$ as an average asymptotic $\bar{\mu}_{123}$-invariant of orbits of $X$, and will become apparent in Section \ref{sec:ergodic}. Formula \eqref{eq:mu_123-helicity} can be also 
 regarded as the third order helicity of three distinct vector fields $X_i$, supported on the handlebodies $\mathcal{T}_i$. In Section \ref{sec:energy} we indicate how to construct the potential $\alpha_\omega$ from the basic elliptic theory of differential forms.

\begin{theorem}[Helicity Invariance Theorem]\label{th:invariance-th}
 On every unlinked invariant handlebody $\mathcal{T}$ in $S^3$, $\mathsf{H}_{123}(X;\mathcal{T})$  is 
 \begin{itemize}
 \item[(i)] independent of a choice of the potential $\alpha_\omega$.
 \item[(ii)] invariant under the action of $\text{Diff}^\text{vol}_0(S^3)$, i.e. for every $g\in \text{Diff}^\text{vol}_0(S^3)$:
\begin{gather}\label{eq:invariance}
 \mathsf{H}_{123}(X;\mathcal{T})=\mathsf{H}_{123}(g_\ast X;g(\mathcal{T}))\ .
\end{gather}
 \end{itemize}
\end{theorem}
\begin{proof}
To prove $(i)$ observe for every other potential $\alpha'_\omega$ of $\omega$, the difference $\beta=\alpha_\omega-\alpha'_\omega$ is a closed $1$-form on $\mathcal{T}$
(since $\omega=d\alpha_\omega=d\alpha'_\omega$). Therefore
\begin{eqnarray*}
 \mathsf{H}_{123}(X;\mathcal{T})-\mathsf{H}_{123}'(X;\mathcal{T}) & = & \int_{\mathcal{T}} (\beta\wedge\omega)\wedge\iota_{X_1}\mu_1\wedge \iota_{X_2}\mu_2\wedge \iota_{X_3}\mu_3\\
 & \stackrel{(1)}{=} & \int_{\mathcal{T}} d\bigl(\beta\wedge \alpha_\omega\wedge\iota_{X_1}\mu_1\wedge \iota_{X_2}\mu_2\wedge \iota_{X_3}\mu_3\bigr)\\
 & = & \int_{\partial\mathcal{T}} \beta\wedge \alpha_\omega\wedge\iota_{X_1}\mu_1\wedge \iota_{X_2}\mu_2\wedge \iota_{X_3}\mu_3 \stackrel{(2)}{=} 0,
\end{eqnarray*}
where in (1) we applied $d(\iota_{X_i}\mu_i)=0$ (since $X_i$'s are divergence free), and 
in (2):
\[
 \iota_{X_i}\mu_i\bigl|_{\partial\mathcal{T}_i}=0,
\]
(because each vector field $X_i$ is tangent to the boundary $\partial \mathcal{T}_i$), and 
\[
 \partial \mathcal{T}=\bigl(\partial \mathcal{T}_1\times \mathcal{T}_2\times \mathcal{T}_3\bigr)\cup \bigl(\mathcal{T}_1\times \partial\mathcal{T}_2\times \mathcal{T}_3\bigr)\cup \bigl(\mathcal{T}_1\times \mathcal{T}_2\times \partial\mathcal{T}_3\bigr).
\]

\no Next we show $(ii)$, the proof is in the style of 
\cite{Berger90, Mayer03}, but adapted to our setting. For any given $g\in \text{Diff}^\text{vol}(S^3)$, there exists a path $t\longrightarrow g(t)\in \text{Diff}^\text{vol}_0(S^3)$,
such that
\[
 g(0)=\text{id}_{S^3},\qquad g(1)=g\ .
\]
 Denote by $V$ the divergence free vector field on $S^3$, given by $V(x)=\frac{d}{dt} g(t,x)|_{t=0}$, i.e. $g(t)$ is a flow of $V$, and push-forward fields $X_i$ by 
\[
 X^t_i:=g(t)_\ast X_i\ .
\]
It is well known (see Appendix B, or \cite[p. 224]{Freedman91-2})
 that 2-forms: $\iota_{X^t_i}\mu$ are frozen in the flow of $V$, i.e.
\begin{equation}\label{eq:forms-frozen-in}
 \frac{d}{dt}\bigl(g(t)^\ast\iota_{X^t_i}\mu\bigr) =(\partial_t+\mathcal{L}_V) \iota_{X^t_i}\mu = 0.
\end{equation} 
We also have a path $\hat{g}(t)=(g(t),g(t),g(t))$ in $\text{Diff}^\text{vol}_0(S^3\times S^3\times S^3)$.
 Analogously for the volume preserving flow $\hat{g}(t)$ we obtain the vector field $\hat{V}=(V,V,V)$ (recall that a tangent bundle $T(S^3)^3$ has a natural product structure). Equation \eqref{eq:forms-frozen-in} implies
\begin{equation}\label{eq:forms-frozen-in2}
(\partial_t+\mathcal{L}_{\hat{V}}) \bigl(\pi^\ast_i\iota_{X^t_i}\mu\bigr)=(\partial_t+\mathcal{L}_{\hat{V}}) \iota_{X^t_i}\mu_i= 0\ .
\end{equation} 
(in the second equation we merely revoke our notational conventions: $\iota_{X^t_i}\mu_i\equiv\pi^\ast_i\bigl(\iota_{g(t)^\ast X_i}\mu\bigr)$).

\no Let $\mathcal{T}(t)=\hat{g}(t)(\mathcal{T}(0))\subset \Cnf(S^3,3)$ we must show $\frac{d}{dt} 
\mathsf{H}_{123}(X^t_1,X^t_2,X^t_3)=0$  at every $t=t_0$. Notice that 
for small enough $\epsilon$ and $t\in (t_0-\epsilon,t_0+\epsilon)$ we can assume, by $(i)$, that $\alpha_\omega$ is a time independent potential obtained from slightly bigger domain $\widetilde{\mathcal{T}}$ which satisfies 
\[
  \mathcal{T}(t)\subset\widetilde{\mathcal{T}},\qquad \forall t\in (t_0-\epsilon,t_0+\epsilon),
\]
and is a deformation retraction of $\mathcal{T}(t_0)$. Without loss of generality we assume 
$t_0=0$, $\hat{g}(0)=\text{id}_{(S^3)^3}$, and calculate at $t_0$:
\begin{eqnarray*}
 \frac{d}{dt} 
 \Bigl(\mathsf{H}_{123}(g(t)^\ast X_1,g(t)^\ast X_2,g(t)^\ast X_3)\Bigr) & = &
 \frac{d}{dt} \int_{\mathcal{T}(t)} \alpha_\omega\wedge d\alpha_\omega\wedge \iota_{X^t_1}\mu_1\wedge \iota_{X^t_2}\mu_2\wedge \iota_{X^t_3}\mu_3\\
 & = & \int_{\mathcal{T}(0)} \frac{d}{dt} \hat{g}(t)^\ast\bigl(\alpha_\omega\wedge d\alpha_\omega\wedge \iota_{X^t_1}\mu_1\wedge \iota_{X^t_2}\mu_2\wedge \iota_{X^t_3}\mu_3\bigr)\\
 & = & \int_{\mathcal{T}(0)} \bigl(\mathcal{L}_{\partial_t+\hat{V}}(\alpha_\omega\wedge d\alpha_\omega)\bigr)\wedge \iota_{X^t_1}\mu_1\wedge \iota_{X^t_2}\mu_2\wedge \iota_{X^t_3}\mu_3,
\end{eqnarray*}
where in the last identity we applied \eqref{eq:forms-frozen-in2} and the product rule for the Lie derivative.
Now because $\omega\wedge\alpha_\omega$ is time independent (for $t\in (t_0-\epsilon,t_0+\epsilon)$) by the Cartan magic formula we obtain
\begin{eqnarray*}
 \mathcal{L}_{\partial_t+\hat{V}}(\alpha_\omega\wedge d\alpha_\omega) & = & \mathcal{L}_{\hat{V}}(\alpha_\omega\wedge d\alpha_\omega) = \iota_{\hat{V}} d(\alpha_\omega\wedge d\alpha_\omega)+ d(\iota_{\hat{V}} (\alpha_\omega\wedge d\alpha_\omega))\\
 & =  & d(\iota_{\hat{V}} (\alpha_\omega\wedge d\alpha_\omega)),
\end{eqnarray*}
since $d(\alpha_\omega\wedge d\alpha_\omega)=\omega\wedge\omega=0$. %Notice that replacing $\omega$ by a different representative of its cohomology class does not influence this calculation.
 Now, the same argument as in the proof of $(i)$ shows that the right hand side of the previous equation is zero, because $X^t_i$ are tangent to the boundary of $\mathcal{T}_i(t)$.
\end{proof}

% \begin{remark}
%  The crucial obstacle in extending this formula to encompass the whole manifold as the invariant set $\mathcal{T}$ is the fact that the potential $\alpha_\omega$ is not globally defined on $\Cnf(S^3,3)$
%  \ldots\ .
% \end{remark}

\begin{remark}
 Notice that the above argument shows that if we replace $\alpha_{\omega}\wedge d\alpha_\omega$ by virtually any closed 3-form $\eta$ on $\Cnf(S^3,3)$ (or $(S^3)^3$) we obtain some invariant under frozen-in-field deformations. Clearly, if $\eta$ is exact we obtain trivial invariants, therefore the only candidates are cohomology classes of $\Cnf(S^3,3)\cong S^3\times S^2$. In dimension $3$, the only candidate cohomology class of $\Cnf(S^3,3)$ is a dual of $S^3$. Based on the considerations in Section \ref{sec:mu-and-hopf}
  one may argue that an invariant obtained this way is trivial. Indeed the cohomology class $\eta$ evaluated on any 3-torus obtained from a 3-component link in $S^3$ via the map $L$ in \eqref{eq:F}
  is zero. Therefore one could apply the approach via the Ergodic Theorem (as in Section \ref{sec:ergodic}) to show that $\int \eta\wedge\iota_{X}\mu_1\wedge \iota_{X}\mu_2\wedge \iota_{X}\mu_3$ 
 defines a trivial invariant.
  The crucial obstacle in extending the formula in \eqref{eq:mu_123-helicity} to encompass the whole $S^3$ is the fact that the potential $\alpha_\omega$ cannot be globally defined on $\Cnf(S^3,3)$.
\end{remark}

\section{The ergodic interpretation of $\mathsf{H}_{123}(X;\mathcal{T})$}\label{sec:ergodic}

The following statement is often seen in literature \cite{Cantarella-DeTurck-Gluck01, Cantarella-DeTurck-Gluck-Teytel00}: 

\begin{quote}
\emph{Helicity measures the extend to which 
vector fields twist and coil around each other.}
\end{quote}

\no A beauty of Arnold's ergodic approach 
to the helicity $\mathsf{H}_{12}(X)$ is that it makes this statement precise, by interpreting 
$\mathsf{H}_{12}(X)$ as an average asymptotic linking number of orbits of $X$. But it also has a
practical application as it allows us to extend our approach to certain invariant sets of $X$. In this section we apply 
this philosophy to our newly defined invariant $\mathsf{H}_{123}(X;\mathcal{T})$, and interpret it
as the average asymptotic $\bar{\mu}_{123}$-invariant of orbits of $X$. 
Moreover this ergodic interpretation leads us to an alternative, more intuitive proof of Helicity Invariance Theorem \ref{th:invariance-th}.

We begin by observing that given a volume preserving vector field $X$ on $M$ and its flow $\Phi_t$, 
we may regard $X$ as three vector fields on $(M)^3$. Therefore $(\Phi,\Phi,\Phi)$ induces a natural $\R^3$ action defined as
\begin{equation}\label{eq:R3-action}
 \mathbf{\Phi}:\R^3\times (M)^3\longrightarrow (M)^3:\qquad \bigl((s,t,u),x,y,z\bigr)\stackrel{\mathbf{\Phi}}{\longrightarrow} (\Phi(s,x),\Phi(t,y), \Phi(u,z)).
\end{equation}
Observe that $\mathbf{\Phi}$ is a volume preserving action on $(M)^3$.  Our analysis is rooted in techniques developed in \cite{Arnold86, Laurence-Avellaneda93, Laurence-Stredulinsky00b}, the main tool is the the following

\begin{theorem}[Multi-parameter Ergodic Theorem, \cite{Becker81}]\label{th:multi-ergodic}
 For any real valued $L^1$-function $F$ on $(M)^3$, the time
averages under the action in \eqref{eq:R3-action}:
\begin{gather*}
\bar{F}(x,y,z)=\lim_{T\to \infty}\frac{1}{T^3}\int^T_0\int^T_0\int^T_0
 F(\Phi(x,s),\Phi(y,t),\Phi(z,u))\,d s\,d t\,d u
\end{gather*}
converge almost everywhere. Moreover, the limit function
$\bar{F}$ satisfies
\begin{itemize}
 \item[(i)] $\|\bar{F}\|_{L^1((M)^3)}\leq \|F\|_{L^1((M)^3)}$,
 \item[(ii)] $\bar{F}$ is invariant under the $\mathbf{\Phi}$-action,
 \item[(iii)] if $(M)^3$ is of finite volume then
\begin{gather}\label{eq:ergodic-integrals-equal}
 \int_{(M)^3}\bar{F} =\int_{(M)^3} F\ .
\end{gather}
\end{itemize}
\end{theorem}

\no We may now define \emph{invariant unlinked domain} $\mathcal{T}$ of $X$ as an arbitrary $\mathbf{\Phi}$-invariant set, with topological closure $\overline{\mathcal{T}}$ which belongs to a larger product of open sets $\widetilde{\mathcal{T}}=\widetilde{\mathcal{T}}_1\times \widetilde{\mathcal{T}}_2\times \widetilde{\mathcal{T}}_3$
in $\Cnf(S^3,3)$, satisfying the following 
\begin{itemize}
\item[(A)] $\widetilde{\mathcal{T}}$ admits a \emph{short path system} $\mathcal{S}$,
\item[(B)] Equation \eqref{eq:omega-exact} holds on $\widetilde{\mathcal{T}}$.
\end{itemize}
\no By a system of short paths on $\widetilde{\mathcal{T}}$, \cite{Vogel03, Khesin98}, we understand 
a collection of  curves $\mathcal{S}=\{\sigma(x,y)\}$ on each open set $\widetilde{\mathcal{T}}_i$ such that 
 
\begin{itemize}
\item[(a)] for every pair of points $x,y \in \widetilde{\mathcal{T}}_i$ there is a curve $\sigma(x,y):I\mapsto \widetilde{\mathcal{T}}_i$ in $\mathcal{S}$, such that $\sigma(0)=x$ and $\sigma(1)=y$,
\item[(b)] the lengths of paths in $\mathcal{S}$ are uniformly 
bounded above by a common constant. 
\end{itemize}


%\no At this point we do not require a stronger notion of short paths since domains $\widetilde{\mathcal{T}}_i$  are disjoint in $(S^3)^3$. 
\no Topologically, every $\mathbf{\Phi}$-invariant 
set is a union of products of orbits of $X$ in $(S^3)^3$, it is often convenient to think of the 
orbits $\mathbf{\Phi}$-action as a foliation of $(S^3)^3$, then $\mathbf{\Phi}$-invariant 
sets are just union of leaves of such foliation.  A fundamental example of an 
invariant unlinked domain is the case of $\mathbf{\Phi}$-invariant 
set $\mathcal{T}$ contained in the product $\widetilde{\mathcal{T}}=\widetilde{\mathcal{T}}_1\times \widetilde{\mathcal{T}}_2\times \widetilde{\mathcal{T}}_3$ of disjoint open unlinked handlebodies $\widetilde{\mathcal{T}}_i$. Notice that we do not require $X$ to be tangent to $\partial\widetilde{\mathcal{T}}_i$ in this case, and $\widetilde{\mathcal{T}}$ always admits a short path system as we describe in the following

\begin{remark}
In \cite{Vogel03}, Vogel shows that on a closed manifold $M$, geodesics always provide a short path system. When $\mathcal{T}$ is contained in the product of unlinked handlebodies $\widetilde{\mathcal{T}}$
we may easily construct such system on $\widetilde{\mathcal{T}}$.
 Because  $\mathcal{\mathcal{T}}_i$ are proper subsets of $S^3$ we generally do not want to use ambient geodesics from $S^3$ as they may not lie entirely 
in $\mathcal{T}_i$. To obtain $\mathcal{S}$ one may always put an artificial Riemannian metric 
on each $\widetilde{\mathcal{T}}_i$ which makes it a complete Riemannian manifold (e.g. by making $\partial\widetilde{\mathcal{T}}_i$ totally geodesic) and choose $\mathcal{S}$ to be geodesics on such Riemannian manifold. We also observe that one may apply a diffeomorphism $g\in \text{Diff}(\widetilde{\mathcal{T}})$ to $\mathcal{S}$ 
which results in the new system $g\mathcal{S}$ on $g(\widetilde{\mathcal{T}})$.
\end{remark}

\no Next theorem shows that the definition of $\mathsf{H}_{123}(X;\mathcal{T})$
in Equation \eqref{eq:helicity} is valid for invariant unlinked domains of $X$ and is an analog of Arnold's Helicity Theorem in our setting

\begin{theorem}[Ergodic interpretation of $\mathsf{H}_{123}(X;\mathcal{T})$]\label{th:ergodic-mu}
 Given $(X;\mathcal{T})$, the following limit (asymptotic $\bar{\mu}_{123}$-invariant of orbits) exists for  almost 
 all $(x,y,z)\in \mathcal{T}$:
\begin{equation}\label{eq:bar-m_X}
 \bar{m}_X(x,y,z)=\lim_{T\to\infty} \frac{1}{T^3}\bar{\mu}_{123}\bigl(\bar{\mathscr{O}}^{X_1}_T(x),\bar{\mathscr{O}}^{X_2}_T(y),\bar{\mathscr{O}}^{X_3}_T(y)\bigr).
\end{equation}
Moreover, 
\begin{equation}\label{eq:H_123-ergodic}
 \mathsf{H}_{123}(X;\mathcal{T})=\int_{\mathcal{T}} \bar{m}_X(x,y,z)\,\mu(x)\wedge\mu(y)\wedge\mu(z)\ .
\end{equation}
\end{theorem}
\begin{proof}
 The proof is similar to the one in e.g. \cite{Vogel03}. Before we start we must point out
 the following identity (valid for any 3-form $\gamma$ on $M\times M\times M$ and vector fields 
 $X_1,X_2,X_3$ on $M$)
 \[
(\iota_{X_1}\iota_{X_2}\iota_{X_3}\gamma)\wedge\mu\wedge\mu\wedge\mu=
\gamma(X_1,X_2,X_3)\,\mu\wedge\mu\wedge\mu=\gamma\wedge\iota_{X_1}\mu\wedge\iota_{X_2}\mu\wedge\iota_{X_3}\mu\ .
 \]
 The first equation follows from the definition, the second equation is a consequence of the fact that 
 $\iota_X$ is an antiderivation i.e.
 \begin{equation}\label{eq:iota-antiderivation}
  \iota_X(\alpha\wedge\beta)=(\iota_X\alpha)\wedge\beta+(-1)^{|\alpha|}\alpha\wedge(\iota_X\beta),
 \end{equation}
   and $\iota_{X_i}\mu_j\equiv 0$ for $i\neq j$. As a result
 \begin{eqnarray*}
\notag  \mathsf{H}_{123}(X_1,X_2,X_3) & = & \int_{\mathcal{T}} \bigl(\alpha_\omega\wedge d\alpha_\omega\bigr)\wedge \iota_{X_1}\mu_1\wedge \iota_{X_2}\mu_2\wedge \iota_{X_3}\mu_3\\
& = & \int_\mathcal{T} \bigl(\iota_{X_1}\iota_{X_2}\iota_{X_3}\bigl(\alpha_\omega\wedge d\alpha_\omega\bigr)\bigr)\,\mu_1\wedge\mu_2\wedge\mu_3\\
\label{eq:m-function}  & = & \int_\mathcal{T} m_{X_1,X_2,X_3}(x,y,z)\,\mu_1\wedge \mu_2\wedge \mu_3,
\end{eqnarray*}
where $m_{X_1,X_2,X_3}:=\alpha_\omega\wedge d\alpha_\omega(X_1,X_2,X_3)$. For convenience let (see \eqref{eq:orbits})
\[
 \bar{\mathscr{O}}(x,y,z;T):=\bar{\mathscr{O}}^{X_1}_T(x)\times \bar{\mathscr{O}}^{X_2}_T(y)\times\bar{\mathscr{O}}^{X_3}_T(y)\ .
\]
\no Clearly, every link $\bar{\mathscr{O}}(x,y,z;T)$ is borromean and by Formula \eqref{eq:mu-formula} we obtain
\begin{eqnarray*}
\bar{\mu}_{123}\bigl(\bar{\mathscr{O}}^{X_1}_T(x),\bar{\mathscr{O}}^{X_2}_T(y),\bar{\mathscr{O}}^{X_3}_T(y)\big)
& = & \int_{\bar{\mathscr{O}}(x,y,z;T)}
\alpha_\omega\wedge d\alpha_\omega\\
& = & \int_{\mathscr{O}(x,y,z;T)}
\alpha_\omega\wedge d\alpha_\omega+(I)
\end{eqnarray*}
where the term $(I)$ involves integrals over short paths in $\mathcal{S}$, (see Appendix B). Next, notice that for every $(x,y,z)$ : $D\mathbf{\Phi}_{(x,y,z)}[\partial_i]=X_i$, $i=s,t,u$ we obtain
\begin{eqnarray*}
 \mathbf{\Phi}^\ast_{(x,y,z)}(\alpha_\omega\wedge d\alpha_\omega) & = & (\alpha_\omega\wedge d\alpha_\omega(X_1,X_2,X_3)(\Phi^1(x,s),\Phi^2(y,t),\Phi^3(z,u))\,ds\wedge dt\wedge du\\
 & = & m_{X_1,X_2,X_3}(\Phi^1(x,s),\Phi^2(y,t),\Phi^3(z,u))\,ds\wedge dt\wedge du\ .
\end{eqnarray*}
Thus 
\[
\int_{\mathscr{O}(x,y,z;T)}
\alpha_\omega\wedge d\alpha_\omega=\int^T_0\int^T_0\int^T_0 m_{X_1,X_2,X_3}(\Phi^1(x,s),\Phi^2(y,t),\Phi^3(z,u))\,ds\wedge dt\wedge du\ .
\]
Function $m_{X_1,X_2,X_3}$ is smooth on $\mathcal{T}$ and therefore $L^1$. Because short paths do not contribute to the time average i.e.
\[
 \lim_{T\to\infty} \frac{1}{T^3}(I)=0
\]
(see Appendix B), Theorem \ref{th:multi-ergodic} applied to the function $m_{X_1,X_2,X_3}$ yields almost everywhere existence of the limit \eqref{eq:bar-m_X} and the invariant $L^1$-function $\bar{m}_X:=\bar{m}_{X_1,X_2,X_3}$ on $\mathcal{T}$. The identity \eqref{eq:H_123-ergodic} 
follows from $(iii)$ of Theorem \ref{th:multi-ergodic}.
\end{proof}

\begin{theorem}[Helicity Invariance Theorem-ergodic version]\label{th:invariance-th-ergodic}
 On every unlinked domain $\mathcal{T}$ in $S^3$, (i) and (ii) of Theorem 
 \ref{th:invariance-th} hold for $\mathsf{H}_{123}(X;\mathcal{T})$.
 \end{theorem}

\begin{proof}%[Proof of $(ii)$ in Theorem \ref{th:invariance-th} (ergodic version)]
The proof of $(i)$ immediately follows from independence of the limit \eqref{eq:bar-m_X}
on the choice of the potential $\alpha_\omega$.
For the proof of  $(ii)$ we must show the following
\[
 \int_\mathcal{T} m_{X_1,X_2,X_3}\,\mu_1\wedge \mu_2\wedge \mu_3=\int_{g(\mathcal{T})} m_{g_\ast X_1,g_\ast X_2, g_\ast X_3}\,\mu_1\wedge \mu_2\wedge \mu_3\ .
\]
Theorem \ref{th:multi-ergodic}, tells us that $m_{X_1,X_2,X_3}$ and $m_{g_\ast X_1,g_\ast X_2, g_\ast X_3}$ admit $L^1$-averages $\bar{m}_{X_1,X_2,X_3}$ and $\bar{m}_{g_\ast X_1,g_\ast X_2, g_\ast X_3}$, under actions of $X_i$ and $g_\ast X_i$ respectively. If the following holds:
\begin{equation}\label{eq:m=gm}
 \bar{m}_{X_1,X_2,X_3}(x,y,z)=\bar{m}_{g_\ast X_1,g_\ast X_2, g_\ast X_3}(g(x),g(y),g(z)),\qquad \text{a.e.},
\end{equation}
then \eqref{eq:invariance} is an immediate consequence of \eqref{eq:ergodic-integrals-equal}, change of variables for integrals, and the fact that $g$ preserves the volume, i.e. $\mu_i=g^\ast\mu_i$.
Borrowing notation from the previous theorem we set
\begin{eqnarray*}
 g\bar{\mathscr{O}}(x,y,z;T) & := & g(\bar{\mathscr{O}}^{X_1}_T(x))\times g(\bar{\mathscr{O}}^{X_2}_T(y))\times g(\bar{\mathscr{O}}^{X_3}_T(z))\\
 & = & \bar{\mathscr{O}}^{g_\ast X_1}_T(g(x))\times \bar{\mathscr{O}}^{g_\ast X_2}_T(g(y))\times \bar{\mathscr{O}}^{g_\ast X_3}_T(g(z)),
\end{eqnarray*}
where the second identity is a consequence of the fact that the flow of $g_\ast X_i$ is 
obtained as a composition of $g$ and the flow of $X_i$.
Since $g$ is isotopic to the identity, $\bar{\mathscr{O}}(x,y,z;T)$ and $\bar{\mathscr{O}}(x,y,z;T)$
 are isotopic borromean links in $S^3$, and for every $T$:
 \[
  \bar{\mu}_{123}(\bar{\mathscr{O}}(x,y,z;T))=\bar{\mu}_{123}(g\bar{\mathscr{O}}(x,y,z;T)).
 \]
 As a result of the above identity and \eqref{eq:bar-m_X} we derive a.e.
\begin{eqnarray*}
  \bar{m}_{X_1,X_2,X_3}(x,y,z) & = & \lim_{T\to\infty}\frac{1}{T^3}\bar{\mu}_{123}(\bar{\mathscr{O}}(x,y,z;T))\\
   & = & \lim_{T\to\infty}\frac{1}{T^3}\bar{\mu}_{123}(g\bar{\mathscr{O}}(x,y,z;T))=\bar{m}_{g_\ast X_1,g_\ast X_2, g_\ast X_3}(g(x),g(y),g(z)).
\end{eqnarray*}
Notice that in the last equation we used in the ``pushed forward'' short paths system: $g\mathcal{S}$.
But since lengths of paths in $g\mathcal{S}$ are bounded as well, they do not contribute to the limit. This proves the identity \eqref{eq:m=gm}, and consequently \eqref{eq:invariance}.
\end{proof}

\no Notice that the above argument does not require the Stokes Theorem, and thus may lead to further generalizations. Clearly, for $\mathsf{H}_{123}(X;\mathcal{T})$ to be nontrivial $\mathcal{T}$ must be 
of nonzero measure.


 The following formula is a well known property of the ordinary helicity $\mathsf{H}_{12}(X;\mathcal{T})$ of the flux tubes $\mathcal{T}$ modeled on a 2-component link $L=\{L_1,L_2\}$ (see e.g. \cite{Laurence-Avellaneda93, Cantarella00})

\begin{equation}\label{eq:helicity12-flux-tubes}
\mathsf{H}_{12}(X;\mathcal{T})=\mathsf{H}_{12}(X_1,X_2)=\text{lk}(L_1,L_2)\,\text{Flux}(X_1)\,\text{Flux}(X_2).
\end{equation}
\no Next we show that an analogous property holds $\mathsf{H}_{123}(X;\mathcal{T})$, for invariant unlinked  handlebodies $\mathcal{T}$ investigated in the previous section.

\begin{proposition}
 $\mathsf{H}_{123}(X_1,X_2,X_3)$ on invariant unlinked handlebodies $\mathcal{T}$ satisfies the following formula
 \begin{equation}\label{eq:helicity-fluxes}
 \mathsf{H}_{123}(X_1,X_2,X_3)=\sum_{i\neq j\neq k} \bar{\mu}_{123}(L^i_1, L^j_2, L^k_3)\,\text{Flux}_{\Sigma_i}(X_1)\,\text{Flux}_{\Sigma_j}(X_2)\,\text{Flux}_{\Sigma_k}(X_3)\ .
 \end{equation}
where $\{L^i_1\otimes L^j_2\otimes L^k_3\}$ is a basis of basic cycles of $H_1(\mathcal{T}_1)\otimes H_1(\mathcal{T}_2)\otimes H_1(\mathcal{T}_3)$, $\text{Flux}(X_i)$ stands for the flux of $X_i$ through a cross sectional surface $\Sigma_k$ of $\mathcal{T}_i$, which represents a homology Poincar$\acute{\rm e}$ dual of $L^k_i$ in $H_2(\mathcal{T}_i,\partial\mathcal{T}_i)$.

\end{proposition}
\begin{proof}
 Recall that a flux $\text{Flux}_{\Sigma_k}(X_i)$ of a vector field $X_i$ though a cross-sectional surface
 $\Sigma_k$ is given by:
 \begin{eqnarray*}
  \text{Flux}_{\Sigma_k}(X_i) & = & \int_{\Sigma_k} \iota_{X_i}\mu=\int_{\mathcal{T}_i} h_k\wedge\iota_{X_i}\mu\\
  & = & \int_{\mathcal{T}_i} \iota_{X_i}h_k\wedge\mu =\int_{\mathcal{T}_i} h_k(X_i),
 \end{eqnarray*}
 where 1-forms $h_k$ represent cohomology Poincar$\acute{\rm e}$ duals of $\Sigma_k$ and we applied \eqref{eq:iota-antiderivation} in the third equation. Clearly, $h_i\in \Omega^1(\mathcal{T}_i)$ satisfy 
 \[
  \int_{\gamma} h_k=\#(\gamma,\Sigma_k)=\text{deg}(\gamma, L_k),
 \]
 where $\text{deg}(\gamma, L^j_k)$ measures how many times $\gamma$ ``wraps around'' in the cycle  $L^j_k$. 
 
 For simplicity we first assume that $\mathcal{T}$ is modeled on a borromean link
 $L=\{L_1,L_2,L_3\}$ (see Figure \ref{fig:invariant-sets}). Then $h=h_1\wedge h_2\wedge h_3\in \Omega^3(\mathcal{T})$  is a cohomology class dual of the cycle $L=L_1\times L_2\times L_3$ in $\mathcal{T}$. We set
 \[
  H:=\iota_{X_1}\iota_{X_2}\iota_{X_3} h=h(X_1,X_2,X_3),
 \]
 which is a smooth function, and let $\bar{H}$ be the time average of $H$ as in Theorem \ref{th:multi-ergodic}. It suffices to show 
 \begin{equation}\label{eq:H-equation}
   \bar{m}_{X_1,X_2,X_3}=\bar{\mu}_{123}(L)\,\bar{H},\qquad \text{a.e.}
 \end{equation}
  Analogously as in Theorem \ref{th:invariance-th}, 
 Equation \eqref{eq:H-equation} immediately implies Formula \eqref{eq:helicity-fluxes}. 
 Assume the notation of Theorem \ref{th:invariance-th}, for given $T$ consider $\bar{\mathscr{O}}(x,y,z;T)$. 
 By the fundamental property of the $\bar{\mu}_{123}$-invariant (see Corollary \ref{th:general-mu_123}) we have 
 \begin{eqnarray*}
  \bar{\mu}_{123}(\bar{\mathscr{O}}(x,y,z;T)) & = & \bar{\mu}_{123}(L)\,\bigl(\text{deg}(\bar{\mathscr{O}}^{X_1}_T(x),L_1)\,\text{deg}(\bar{\mathscr{O}}^{X_2}_T(y),L_2)\,\text{deg}(\bar{\mathscr{O}}^{X_3}_T(x),L_3)\bigr)\\
  & = & \bar{\mu}_{123}(L)\int_{\bar{\mathscr{O}}(x,y,z;T)} h\ .
 \end{eqnarray*}
Thus we obtain 
 \begin{eqnarray*}
  \bar{m}_{X_1,X_2, X_3}(x,y,z) & = & \lim_{T\to \infty} \frac{1}{T^3}\bar{\mu}_{123}(\bar{\mathscr{O}}(x,y,z;T))\\
  & = & \bar{\mu}_{123}(L)\lim_{T\to \infty} \frac{1}{T^3}\int_{\bar{\mathscr{O}}(x,y,z;T)} H\\
  & =  &\bar{\mu}_{123}(L) \bar{H}(x,y,z),
 \end{eqnarray*}
where the last equality is again the consequence of short paths not contributing to the limit.

\no Now, from the product structure and $(iii)$ of Theorem \ref{th:multi-ergodic} we have
 \[
  \int_\mathcal{T} \bar{H} =\int_\mathcal{T} H=\bigl(\int_{\mathcal{T}_1} h_1(X_1) \bigr)\bigl(\int_{\mathcal{T}_2} h_2(X_2) \bigr)\bigl(\int_{\mathcal{T}_3} h_3(X_3) \bigr).
 \]
The proof in the case of a general handlebody is analogous, once we show the following 
\begin{lemma}
 Let $\mathcal{O}=\{\mathcal{O}_1,\mathcal{O}_2, \mathcal{O}_3\}$ be a 3-component link in $S^3$ such that 
 $\mathcal{O}_i\subset \mathcal{T}_i$ for each $i$. Then $\mathcal{O}$ is borromean and 
  \begin{equation}\label{eq:mu-in-T}
 \bar{\mu}_{123}(\mathcal{O})=\sum_{i\neq j\neq k} \bar{\mu}_{123}(L^i_1, L^j_2, L^k_3)\,\text{deg}(\mathcal{O}_1, L^i_1)\,\text{deg}(\mathcal{O}_2, L^j_2)\,\text{deg}(\mathcal{O}_3, L^k_1)\ .
 \end{equation}
\end{lemma}
\begin{proof}
 Notice that thanks to our new interpretation, the $\bar{\mu}_{123}$-invariant is not only a link homotopy 
 invariant but a homotopy invariant of the associated map $F_{\mathcal{O}}$ defined in \eqref{eq:F}. Observe that  each component 
 $\mathcal{O}_i$ can be homotopied inside of its handlebody $\mathcal{T}_i$ to be come a bouquet of circles $\widehat{\mathcal{O}}_i\cong S^1\vee S^1\vee\ldots\vee S^1$. Such that each factor in $\widehat{\mathcal{O}}_i$ is a basic cycle in $H_1(\mathcal{T}_i)$. As a result we obtain the associated map
 $F_{\widehat{\mathcal{O}}}$, and
 \[
  \bar{\mu}_{123}(\mathcal{O})=\frac{1}{2}\mathscr{H}(F_{\widehat{\mathcal{O}}}).
 \]
 Interpreting, $\mathscr{H}(F_{\widehat{\mathcal{O}}})$ as the intersection number and summing up
 intersection numbers we conclude \eqref{eq:mu-in-T}.
\end{proof}

\end{proof}

\no In the case of borromean flux tubes modeled on $L$, \eqref{eq:helicity12-flux-tubes} reduces to
\begin{equation}\label{eq:helicity12-borro-flux-tubes}
\mathsf{H}_{123}(X_1,X_2,X_3)=\bar{\mu}_{123}(L)\,\text{Flux}_{\Sigma_1}(X_1)
\text{Flux}_{\Sigma_2}(X_2)\text{Flux}_{\Sigma_3}(X_3),
\end{equation}
where $\Sigma_i$ denotes homology Poincar$\acute{\rm e}$ duals to $L_i$ in $H_2(\mathcal{T}_i,\partial\mathcal{T}_i)$.


  Since the fluxes are invariant under frozen-in-field deformations Formula \eqref{eq:helicity12-flux-tubes}
  is yet another proof of Theorem \ref{th:invariance-th} in the setting of invariant unlinked handlebodies.
  In \cite{Laurence-Stredulinsky00b} the authors develop the same formula for the borromean flux tubes.
  This clearly must be the case as we deal with the same topological 
  invariants of links via a different approach. Additional advantage 
  of our formulation is that we do not have to separately 
  deal with null points of vector fields as in \cite{Laurence-Stredulinsky00b}. We also note that 
  in  \cite{Laurence-Stredulinsky00b} the authors prove the helicity invariance only via Formula \eqref{eq:helicity12-borro-flux-tubes}.



\section{Energy bound.}\label{sec:energy}

\no In this section we indicate how the quantity $\mathsf{H}_{123}(X;\mathcal{T})$ invariant under frozen-in-field deformations provides a lower bound for the  $L^2$-energy $E_2(X)$ of a volume preserving field $X$ on $M=S^3$. For simplicity we restrict our considerations to the case of an invariant unlinked  handlebody  $\mathcal{T}$, as in Section \ref{sec:handlebody}. For the notation used in this section see Appendix C.


Recall the definition
\begin{equation}\label{eq:energy}
 E_2(X)=\int_M |X|^2 = \|X\|^2_{L^2(M)}\ .
\end{equation}
The ordinary helicity $\mathsf{H}_{12}(X)$ provides the well known lower bound (see \cite[p. 123]{Khesin98}):
\[
  \frac{1}{\lambda_1}|\mathsf{H}_{12}(X)|\leq  E_2(X),
\]
where $\lambda_1$ is the first eigenvalue of the elliptic self adjoint operator $\ast\,d:\Omega^1(M)\mapsto \Omega^1(M)$
(i.e. the curl operator). Importance of such lower energy bounds,
% in terms of invariants of volume preserving flows defined in \eqref{eq:invariants},
 stems from an area of interest in the ideal magneto-hydrodynamics, \cite{Priest84}, as this constrains the phenomenon of ``magnetic relaxation''.
A need for higher helicities can be justified by the fact that one may easily produce examples of vector fields $X$ for which $\mathsf{H}_{12}(X)$ vanishes, but the energy of the field $X$ still cannot be relaxed. For example consider a classical case of borromean flux tubes $\mathcal{T}^\text{Borr}$ with $X$ smooth and vanishing outside the tubes. Moreover assume that orbits of $X$ are just ``parallel 
circles'' inside each tube. By taking $X=X_1+X_2+X_3$, bilinearity of $\mathsf{H}_{12}(\,\cdot\,,\,\cdot\,)$ on $\text{Vect}^\text{vol}(S^3)$ and disjoint supports of $X_i$  yield
\[
 \mathsf{H}_{12}(X)=\mathsf{H}_{12}(X,X)=\sum^3_{i=1}\mathsf{H}_{12}(X_i,X_i)+\sum_{i\neq j} \mathsf{H}_{12}(X_i,X_j)=0,
\]
because ``cross-helicities'' $\mathsf{H}_{12}(X_i,X_j)$, $i\neq j$ vanish by Formula \eqref{eq:helicity12-flux-tubes}, and self-helicities $\mathsf{H}_{12}(X_i,X_i)$ vanish as well because the average linking number of orbits is zero (orbits are just parallel 
circles). Nevertheless, Formula \eqref{eq:helicity-fluxes} tells us 
\[
\mathsf{H}_{123}(X;\mathcal{T}^\text{Borr})\neq 0,
\]
thus we may regard $\mathsf{H}_{123}(X;\mathcal{T})$ as a possible ``higher obstruction'' to the energy relaxation or the third order cross-helicity of $X$ on $\mathcal{T}$.


%we will also consider a slightly bigger domain $\widetilde{\mathcal{T}}$, containing $\mathcal{T}$, with smooth boundary, which is a tubular neighborhood of $\mathcal{T}$ in $\Cnf(S^3,3)$. 

%Next, we consider slightly bigger disjoint open sets $U_j$ so that the domain $U=U_1\,\cup\, U_2\,\cup\, U_3$ deformation retracts onto $\mathcal{T}$. Let $U=U_1\times U_2\times U_3\subset \Cnf(S^3,3)$, w.l.o.g. we may assume that $\partial U$  is smooth.
%Now from the product structure the homology $H_1(U)$ is generated by cycles: $L_j$ (appropriately embedded in $U$),
%and $H_2(U)$ by cycles $L_{ij}=L_i\times L_j$. By duality the cohomology $\mathcal{H}^2_N(U)$ is generated by 
%closed 2-forms $h_{12}, h_{23}, h_{13}$ satisfying the Neumann's boundary conditions and such that 
%\begin{equation}\label{eq:h_kl-dual}
%   \int_{L_{ij}} h_{kl} =\delta_{ij,kl}\ .
%\end{equation}
%Consider the closed form $\omega=H^\ast(\nu)\in \Omega^2(\Cnf(S^3,3))$, clearly $\omega$ is well defined smooth form on $U$ since $U\subset \Cnf(S^3,3)$. Assume that $U$ has a product Riemannian metric induced from $S^3\times S^3\times S^3$. By the Fredrichs decomposition we have
%\[
% \omega=l_{12}\,h_{12}+l_{23}\,h_{23}+l_{13}\,h_{13}+d\alpha_\omega=d\alpha_\omega,
%\]
%where the second equality follows from the fact that $l_{ij}=\text{lk}(L_i,L_j)=0$ (by assumption that $L$ is borromean). Indeed 
%\begin{eqnarray}
%\label{eq:condition1} \text{lk}(L_i,L_j) & = & \int_{L_i\times L_j} \omega = \int_{L_i\times L_j} l_{ij} h_{ij} +\int_{L_i\times L_j} d\alpha_\omega\\
%\notag  & = & l_{ij}=0,
%\end{eqnarray}
%for some constants $l_{ij}$.
%
%
%Recall that the $L^2$-energy $E(X)$ of a vector field $X$ on $M$ is defined as
%
%
%Given a vector field $X$ on $M=S^3$ assume that there exist an invariant domain $\mathcal{T}=\mathcal{T}_1\cup \mathcal{T}_2\cup \mathcal{T}_3\subset S^3$ of $X$ in the ``shape'' of pairwise unlinked disjoint borromean tubes (such that $X$ is tangent to $\partial \mathcal{T}_i$), and denote by $X_i=X|_{\partial \mathcal{T}_i}$, $i=1,2,3$. Clearly, in this situation the pairwise helicities 
%$i,j=1,2,3$, $i\neq j$, $\mathscr{H}(X_i,X_j)=0$ vanish so they do not provide a bound for $E(X)$. The self-helicities
%$\mathscr{H}(X_i,X_i)$ may provide such a bound\footnote{for nice ``parallel'' field lines this is not the case however}, depending on the topology of the field lines inside the tubes, but our purpose here is to show that \emph{merely existence of such borromean invariant tori} (i.e. three invariant toroidal surfaces $\partial \mathcal{T}_i$) \emph{leads to a bound for $E(X)$, if the $\bar{\mu}_{123}$-invariant of $\mathcal{T}$ is nonzero.}
%

 To obtain a lower bound for $E_2(X)$ in such situations we notice that 
$\mathsf{H}_{123}(X;\mathcal{T})$ is the $L^2$-inner product of the $6$-forms:
$\ast(\omega\wedge\alpha_\omega)$ and $\iota_{X_1}\mu_1\wedge \iota_{X_2}\mu_2\wedge \iota_{X_3}\mu_3$.
Indeed from \eqref{eq:mu_123-helicity} we have
\[
 \mathsf{H}_{123}(X;\mathcal{T})=\bigl(\ast(\alpha_\omega\wedge \omega),\iota_{X_1}\mu_1\wedge \iota_{X_2}\mu_2\wedge \iota_{X_3}\mu_3\bigr)_{L^2\Omega^6((S^3)^3)}\ .
\]
Let $C_\mathcal{T}=\|\alpha_\omega\wedge \omega\|_{L^2}$, for a fixed Riemannian product metric on $(S^3)^3$ this constant depends on the domain $\mathcal{T}$ in $(S^3)^3$. We estimate using the Cauchy-Schwartz inequality
\begin{eqnarray}
\notag |\mathsf{H}_{123}(X;\mathcal{T})| & \leq & \|\alpha_\omega\wedge \omega\|_{L^2}\|\iota_{X_1}\mu_1\wedge \iota_{X_2}\mu_2\wedge \iota_{X_3}\mu_3\|_{L^2}\\
\label{eq:energy-est1} & = & C_\mathcal{T}\Bigl(\int_{\mathcal{T}}(\iota_{X_1}\mu_1\wedge \iota_{X_2}\mu_2\wedge \iota_{X_3}\mu_3)\wedge \ast(\iota_{X_1}\mu_1\wedge \iota_{X_2}\mu_2\wedge \iota_{X_3}\mu_3)\Bigr)^\frac{1}{2}\\
\notag & \stackrel{(1)}{=} & C_\mathcal{T}\Bigl(\int_{\mathcal{T}}|X_1|^2|X_2|^2|X_3|^2\,\mu\Bigr)^\frac{1}{2}\\
\notag & \stackrel{(2)}{=} & C_\mathcal{T}\Bigl(\int_{\mathcal{T}_1}|X_1|^2\Bigr)^\frac{1}{2}\Bigl(\int_{\mathcal{T}_2}|X_2|^2\Bigr)^\frac{1}{2}\Bigl(\int_{\mathcal{T}_3}|X_3|^2\Bigr)^\frac{1}{2}\\
\label{eq:energy-est} & \leq &   C_\mathcal{T} E_2(X)^\frac{3}{2},
\end{eqnarray}
where $E_2(X)=\int_{S^3} |X|^2$.
To observe (1) first note that for any pair of forms: $\alpha\in \Omega^k(M)$, and $\beta\in \Omega^j(N)$, on Riemannian manifolds $M$ and $N$, on the product $M\times N$ we have  \begin{equation}\label{eq:hodge-product}
  \ast_{M\times N}(\pi^\ast_M(\alpha)\wedge \pi^\ast_N(\beta))=(\ast_M \pi^\ast_M \alpha)\wedge (\ast_N \pi^\ast_N \beta),
 \end{equation}
where $\pi_M:M\times N\longrightarrow M$ and $\pi_N:M\times N\longrightarrow N$ are the natural projections, (the proof is a simple calculation in an orthogonal frame of the product and is left to the reader).
Now, identity (1) follows by applying \eqref{eq:hodge-product} to the integrand in \eqref{eq:energy-est1}, and observing in the coframe $\{\eta^i_k\}$ on the $i$-th factor of $(S^3)^3$
 ($X_i=(a_1,a_2,a_3)$):
\begin{eqnarray*}
 \iota_{X_i}\mu_i\wedge \ast \iota_{X_i}\mu_i & = & (a_1\eta^i_2\wedge\eta^i_3-a_2\eta^i_1\wedge\eta^i_3+a_3\eta^i_1\wedge\eta^i_2)\wedge (a_1\eta^i_1-a_2\eta^i_2+a_3\eta^i_3)\\
 & = & (a^2_1+a^2_2+a^2_3)\, \eta^i_1\wedge\eta^i_2\wedge\eta^i_3 = |X_i |^2 \eta^i_1\wedge\eta^i_2\wedge\eta^i_3=|X_i |^2 \mu_i\ .
\end{eqnarray*}
The next identity (2) follows from the Fubini's Theorem.

Our next goal is to provide an estimate for $C_\mathcal{T}$. For this purpose we review 
some basic $L^2$-theory of the operator $d^{-1}$ (i.e. inverse of the exterior derivative $d$).
The main goal is to estimate an $L^2$-norm of the potential $\alpha_\omega$ of $\omega$ in \eqref{eq:omega-exact}.


Following the standard elliptic theory of differential forms, \cite{Schwarz95}, the potential $\alpha_\omega$ in \eqref{eq:omega-exact} can be obtained via a solution to the Neumann problem for $2$-forms on $\widetilde{\mathcal{T}}$ (see Appendix C)
\begin{equation}\label{eq:neumann-problem}
 \begin{cases}
  \Delta \phi_N=\omega,\qquad \text{in } \widetilde{\mathcal{T}}\\
  \mathbf{n}\,\phi_N=\mathbf{n}\,d\phi_N=0,\qquad \text{on } \partial \widetilde{\mathcal{T}},
 \end{cases}
\end{equation}
where $\mathbf{n}$ stands for the normal component of a differential form, and $\Delta=d\delta+\delta d$.
Because $\mathcal{T}$
is a domain with corners we replace it by a slightly larger domain $\widetilde{\mathcal{T}}$ in $(S^3)^3$ with the same topology 
(i.e. $\mathcal{T}$ is a deformation retract of $\widetilde{\mathcal{T}}$) but with smooth boundary $\partial\widetilde{\mathcal{T}}$.
Notice that because of the ``gauge invariance'' $(i)$ of Theorem \ref{th:invariance-th}, in the definition of $\mathsf{H}_{123}(X;\mathcal{T})$ we my use the restriction of 
\begin{equation}\label{eq:potential-neumann}
\alpha_\omega=\delta\phi_N,
\end{equation}
 to $\mathcal{T}$ (see Appendix C for justification of \eqref{eq:potential-neumann}).
Associated to \eqref{eq:neumann-problem} is the Neumann Laplacian
\[
 \Delta_N:H^2\Omega^2_N(\widetilde{\mathcal{T}})\longrightarrow L^2\Omega^2(\widetilde{\mathcal{T}}),
\] 
which has a discrete positive spectrum $\{\lambda_{i,N}\}$ and eigenvalues satisfy the variational principle called \emph{Rayleigh-Ritz quotient}, \cite{Chavel84}. The first (principal) eigenvalue  $\lambda_{1,N}$ may be expressed as
\[
 \lambda_{1,N}=\inf\left( 	\frac{\|d\varphi\|^2_{L^2}+\|\delta\varphi\|^2_{L^2}}{\|\varphi\|^2_{L^2}}\ \Bigl|\ \varphi\in H^1\Omega^2_N(M)\cap \mathcal{H}^2_N(M)^\perp\right)
\]
\no We denote the inverse of $\Delta_N$ by
\[
 G_N:L^2\Omega^2(\widetilde{\mathcal{T}})\longrightarrow H^2\Omega^2_N(\widetilde{\mathcal{T}}),
\]
which restricts to a compact, self-adjoint operator on $L^2$. As a result spectrum of $G_N$ is 
discrete and given as $\{1/\lambda_{i,N}\}$. Note that based on these considerations we may define $d^{-1}:=\delta G_N$.


\begin{theorem}[Energy bound]
 For every volume preserving vector field $X$ which has an invariant unlinked domain $\mathcal{T}$,
 the $L^2$-energy of $X$ on $S^3$ is bounded below by the third order helicity $\mathsf{H}_{123}(X;\mathcal{T})$, as follows
 \begin{equation}
  |\mathsf{H}_{123}(X;\mathcal{T})|\leq C_\mathcal{T} (E_2(X))^\frac{3}{2}.
 \end{equation}
Moreover, we may estimate the constant $C_\mathcal{T}$:
 \begin{equation}\label{eq:C_T-estimate}
  C_\mathcal{T}\leq \frac{1}{\sqrt{\lambda_{1,N}}}\|\omega\|^2_{L^\infty\Omega^2(\widetilde{\mathcal{T}})}(\text{Vol}(S^3))^3,
 \end{equation}
where $\lambda_{1,N}$ is the first eigenvalue of the Neumann  Laplacian on  $\Omega^2(\widetilde{\mathcal{T}})$.
\end{theorem}
\begin{proof}
We estimate
\begin{eqnarray*}
 \|\alpha_\omega\|^2_{L^2} = \|\delta \phi_N\|^2_{L^2} & \leq & \|d\phi_N\|^2_{L^2}+\|\delta\phi_N\|^2_{L^2}\\
 & = & |(\Delta \phi_N,\phi_N)| \leq  \|\omega\|_{L^2}\|\phi_N\|_{L^2},
\end{eqnarray*}
where we used the Green's formula \cite[p. 60]{Schwarz95} and boundary conditions of \eqref{eq:neumann-problem} in the second identity. Now, because $\phi_N=G_N\omega$, and 
it is a well known fact that $\|G_N\|_{L^2}=\frac{1}{\lambda_{1,N}}$, ($G_N$ is compact self-adjoint on $L^2$) we obtain:
\[
\|\alpha_\omega\|_{L^2}\leq \frac{1}{\sqrt{\lambda_{1, N}}} \|\omega\|_{L^2}\ .
\]
As a result we estimate $C_\mathcal{T}$:
\begin{eqnarray*}
 C_\mathcal{T}=\|\alpha_\omega\wedge\omega\|_{L^2} & \leq & \|\omega\|_{L^\infty}\|\alpha_\omega\|_{L^2}\\
 & \leq &  \frac{1}{\sqrt{\lambda_{1, N}}} \|\omega\|_{L^2}\|\omega\|_{L^\infty}.
\end{eqnarray*}
%thus \eqref{eq:C_T-estimate} follows.
\end{proof}
Notably the best energy estimate has been obtained by Freedman and He, \cite{Freedman91-2}, for the $L^{3/2}$-energy of $X$,  in the case when $X$ admits an invariant domain $\mathcal{T}$ modeled on an $n$-component link $L=\{L_1,\ldots,L_n\}$ in $\R^3$. Their estimate
is based on the asymptotic crossing number and reads
\begin{equation}\label{eq:energy-free-he}
 E_{3/2}(X)\geq \Bigl(\frac{\pi}{16}\Bigr)^{1/4}\left(\sum^n_{k=1} \text{ac}(L_k,L)|\text{Flux}(X_k)|\right)\cdot \min_{1\leq k\leq n}\{|\text{Flux}(X_k)|\},
\end{equation}
where the asymptotic crossing numbers $\text{ac}(L_k,L)$  for borromean links can be estimated below by
a smallest genus among surfaces in $\R^3\setminus \{L_1\cup\ldots\cup\hat{L}_k\cup\ldots\cup L_n\}$
with a single boundary component $L_i$.  Since $L^{3/2}$-energy of $X$ bounds the $L^2$-energy, via the H$\ddot{\rm o}$lder Inequality, thus 
\eqref{eq:energy-free-he} leads to a lower estimate purely in terms of fluxes and topological data.
It is not clear to the author if this approach can be extended to the case of invariant handlebodies 
considered in Section \ref{sec:handlebody}.
A different, more optimal estimate, has been obtained by Laurence and Stredulinsky, via the Massey product formula, in \cite{Laurence-Stredulinsky00a}, but the proof is provided only in a special case of the vector field $X$.

Contrary to these lower bounds which are given in terms of topological data the estimate in \eqref{eq:C_T-estimate} depends on the geometry of the domain $\mathcal{T}$, and also $\|\omega\|_{L^\infty}$. Unfortunately, $\omega$ blows up on the diagonals $\mathbf{\Delta}\subset (S^3)^3$, and as a result the estimate is meaningless when the handlebodies 
$\mathcal{T}_i$ get close to each other during the evolution of the magnetic field $X$. At this point, we need an assumption for $\mathcal{T}_i$ to stay \emph{1cm} apart during the evolution.
Another drawback is that $\lambda_{1,N}$ is a geometric constant which is altered during the evolution as well. 





\section{Comparison to the known approaches via Massey products.}\label{sec:massey}

In several prior works \cite{Berger90, Berger-Evans92, Laurence-Stredulinsky00b, Mayer03} higher helicities were developed via the Massey product formula for $\bar{\mu}_{123}$. These approaches should be equivalent to the one presented here in the sense that invariants obtained this way measure the same topological information. Most notably the work \cite{Laurence-Stredulinsky00b} provides an explicit expression for the third order helicity of the borromean flux tubes, where the ergodic interpretation in the style of Arnold's asymptotic linking number is also provided.  In  \cite{Mayer03} one finds the following formula for the third order helicity
\begin{equation}\label{eq:massey-3A}
 \mathsf{M}_{123}(X_1,X_2,X_3)=\int_M A_1\wedge A_2\wedge A_3,
\end{equation}
where $A_i=d^{-1}(\iota_{X_i}\mu)$. This formula is valid for three distinct vector fields $X_i$
on a closed manifold $M$. For invariant domains with boundary 
\eqref{eq:massey-3A} defines an invariant provided $A_i\bigl|_{\partial M}=0$, but this only happens in certain situations (e.g. $M$ is simply connected, and $A_i$'s are appropriately chosen). 

\no The most commonly known formula directly related to the Massey products was developed by Berger \cite{Berger90}, in the case of 
borromean flux tubes: $\mathcal{T}=\mathcal{T}_1\cup \mathcal{T}_2\cup \mathcal{T}_3$:
\begin{equation}\label{eq:massey123}
 \mathsf{M}_{123}(X_1,X_2,X_3)=\int_{\partial \mathcal{T}_1} A_1\wedge F_{23}+F_{12}\wedge A_3,
\end{equation}
in \cite{Berger90} it is expressed as a volume integral over $\mathcal{T}$ by applying so called gauge fixing. When $\mathcal{T}_i$ are topologically solid tori there exists a single Massey product $<a_1,a_2,a_3>$ in the complement: $S^3\setminus \mathcal{T}$, and the 2-form $A_1\wedge F_{23}+F_{12}\wedge A_3$ is its representative. When $\mathcal{T}_i$ are handlebodies there are multiple Massey products  in the complement $S^3\setminus \mathcal{T}$, but the formula 
\eqref{eq:massey123} should still be valid. Such extensions were not so far considered in the literature and the volume integrals over $\mathcal{T}$ may be harder to obtain in such case.
One may also point out that ergodic interpretations 
of Massey products are more involved, \cite{Laurence-Stredulinsky00b}, comparing to the approach presented in Section \ref{sec:ergodic}. 



\section*{Appendix A: Equations for the \emph{frozen-in-field} forms}
Given a volume preserving vector field $X$ on $M$ and a path $t\longrightarrow g(t)\in \text{Diff}^\text{vol}(M)$, let 
\[
 X^t:=g_\ast(t)X\ . 
\]
Then for $X^0=X$, by definition 
\begin{equation}\label{eq:euler2}
 \dot{X^t}=\frac{d}{dt} X^t\bigl|_{t=0}=-\Lie_V X=-[V,X],
\end{equation}
and as a result 
\[
  \frac{d}{dt} \iota_{X^t}\mu|_{t=0}  =  -\iota_{[V,X]} \mu  
\]
Next, calculate ($\partial_t+\Lie_V\equiv \Lie_{\partial_t+V}$)
\begin{eqnarray*}
 (\partial_t+\Lie_V)\iota_{X^t}\mu & = & \partial_t(\iota_{X^t}\mu)+\Lie_V\iota_{X^t}\mu\\
 & = & \iota_{\dot{X^t}}\mu+\iota_{X^t}(\Lie_V\mu)+\iota_{[V,X^t]}\mu\\
 & = & \iota_{\dot{X^t}+[V,X^t]}\mu,
\end{eqnarray*}
where in the second identity we applied the general formula: $\iota_{[A,B]}=\Lie_A\iota_B-\iota_B\Lie_A$, and in the third equation the fact that $V$ is volume preserving i.e. $\Lie_V\mu=0$. As a result 
of \eqref{eq:euler2} we obtain
\[
 (\partial_t+\Lie_V)\iota_{X^t}\mu=0\ .
\]

\section*{Appendix B: Zero contribution of short paths to the time average}

It is clear when $\beta\in \Omega^3(\widetilde{\mathcal{T}})$ is at least a $C^1$ on
$\widetilde{\mathcal{T}}\subset S^3\times S^3\times S^3$, then
$f=\beta(X_1,X_2,X_3)$ is continuous on $\widetilde{\mathcal{T}}$ and
\begin{gather*}
\int_{\bar{\mathscr{O}}^1_T(x)}\int_{\bar{\mathscr{O}}^2_T(y)}\int_{\bar{\mathscr{O}}^3_T(z)}
f=\\
\bigl(\int_{\mathscr{O}^1_T(x)}+\int_{\sigma(\Phi^1(x,T),x)}\bigr)
\bigl(\int_{\mathscr{O}^2_T(y)}+\int_{\sigma(\Phi^2(y,T),y)}\bigr)
\bigl(\int_{\mathscr{O}^3_T(z)}+\int_{\sigma(\Phi^3(z,T),z)}\bigr)f\ .
\end{gather*}
After expanding, it is obvious that we must show the following (for all choices), when $T\to \infty$:
\begin{eqnarray}\label{eq:lim-sp1}
\frac{1}{T^3}\int_{\sigma(\Phi^1(x,T),x)}\int_{\mathscr{O}^2_{T}(y)}\int_{\mathscr{O}^3_{T}(z)} f & \longrightarrow & 0,\\
\label{eq:lim-sp2}\frac{1}{T^3}\int_{\sigma(\Phi^1(x,T),x)}\int_{\sigma(\Phi^2(y,T),y)}\int_{\mathscr{O}^3_{T}(z)} f & \longrightarrow & 0,\\
\label{eq:lim-sp3}\frac{1}{T^3}\int_{\sigma(\Phi^1(x,T),x)}\int_{\sigma(\Phi^2(y,T),y)}\int_{\sigma(\Phi^3(z,T),z)} f & \longrightarrow & 0.
\end{eqnarray}
Since the lengths of the short paths in $\mathcal{S}$ are bounded by a common constant $d$, \eqref{eq:lim-sp1}-\eqref{eq:lim-sp3} follow immediately, e.g. for \eqref{eq:lim-sp1} we have

\begin{gather*}
\bigl|\frac{1}{T^3}\int_{\sigma(\Phi^1(x,T),x)}\int_{\mathscr{O}^2_{T}(y)}\int_{\mathscr{O}^3_{T}(z)} f\bigr|\leq \frac{1}{T^3} d(T+d)(T+d)\|f\|_{\infty} \longrightarrow  0.
\end{gather*} 


\section*{Appendix C: Notation}
 We adopt notation from the elegant exposition in \cite{Schwarz95}. Let $M$ be an orientable manifold with smooth boundary.
\begin{itemize}
\item[] $\Omega^k(M)=C^\infty(M,\Lambda^k)$, smooth differential forms on $M$.
\item[] $\Omega^k_N(M)=\{\phi\in \Omega^k(M)\ |\ \mathbf{n}\phi=0,\, \mathbf{n}d\phi=0\}$ the subspace satisfying the Neumann boundary conditions, ($\mathbf{n}$ denotes a normal component of a form along $\partial M$).
\end{itemize}
\no The $L^2$-inner product on $\Omega^k(M)$ is defined as 
 \[
  \bigl( \omega,\eta\bigr)_{L^2}=\int_M \omega\wedge\ast \eta,
 \] 
\begin{itemize}
\item[] $L^2\Omega^k(M)$, $L^2$-differential forms on $M$.
\item[] $H^s\Omega^k(M)$ the Sobolev space of differential $k$-forms with $s$ weak ``derivatives'' in $L^2$.
\item[] $\mathcal{H}^k_N(M)=\{\lambda\in H^1\Omega^k(M)\ |\ d\lambda=\delta\lambda=0, \mathbf{n}\lambda=0\}$ the subspace of the Neumann harmonic fields.
\end{itemize}


\no Next we justify \eqref{eq:potential-neumann}, first observe that for any $\gamma\in H^1\Omega^{k-1}(M)$:
\begin{equation}\label{eq:perp-neumann}
 (d\gamma,\lambda)_{L^2}=\int_{\partial M} \mathbf{t}\gamma\wedge\ast\mathbf{n}\lambda=0,\qquad \forall \lambda\in \mathcal{H}^k_N(M),
\end{equation}
as a result if $\omega \in \mathcal{H}^k_N(M)^\perp$ we obtain a solution $\phi$ to the Neumann problem:
\begin{equation}\label{eq:neumann-phi}
 \delta d\phi+d\delta\phi=\omega,\quad \Rightarrow\quad \omega-d\delta\phi=\delta d\phi\ .
\end{equation}
Formula \eqref{eq:perp-neumann} implies: $(\omega-d\delta\phi)\in \mathcal{H}^k_N(M)^\perp$, moreover
$\mathbf{n}(\omega-d\delta\phi)=\mathbf{n}(\delta d\phi)=\delta\,\mathbf{n}(d\phi)=0$, by the boundary condition in \eqref{eq:neumann-problem}. If $\omega$ is a closed form, $\omega-d\delta\phi$ is also closed,
and clearly coclosed by \eqref{eq:neumann-phi}. Thus $\omega-d\delta\phi$ is a harmonic field with zero normal component, and therefore it has to be in $\mathcal{H}^k_N(M)$, and therefore the zero form. This yields 
\[
 \omega=d\delta\phi\ .
\]
As a result we obtain a necessary and sufficient conditions for $\omega$ to be exact:
 \begin{itemize}
 \item[$(i)$] $d\omega=0$,
 \item[$(ii)$] $(\omega,\lambda)_{L^2}=0$, $\forall\lambda\in \mathcal{H}^k_N(M)$.
\end{itemize}


\begin{thebibliography}{10}

\bibitem{Akhmetiev05}
P.~Akhmetiev.
\newblock On a new integral formula for an invariant of 3-component oriented
  links.
\newblock {\em J. Geom. Phys.}, 53(2):180--196, 2005.

\bibitem{Arnold86}
V.~Arnold.
\newblock The asymptotic {H}opf invariant and its applications.
\newblock {\em Selecta Math. Soviet.}, 5(4):327--345, 1986.
\newblock Selected translations.

\bibitem{Khesin98}
V.~Arnold and B.~Khesin.
\newblock {\em Topological methods in hydrodynamics}, volume 125 of {\em
  Applied Mathematical Sciences}.
\newblock Springer-Verlag, New York, 1998.

\bibitem{Becker81}
M.~E. Becker.
\newblock Multiparameter groups of measure-preserving transformations: a simple
  proof of {W}iener's ergodic theorem.
\newblock {\em Ann. Probab.}, 9(3):504--509, 1981.

\bibitem{Berger90}
M.~Berger.
\newblock Third-order link integrals.
\newblock {\em J. Phys. A}, 23(13):2787--2793, 1990.

\bibitem{Bott82}
R.~Bott and L.~Tu.
\newblock {\em Differential forms in algebraic topology}, volume~82 of {\em
  Graduate Texts in Mathematics}.
\newblock Springer-Verlag, New York, 1982.

\bibitem{Cantarella00}
J.~Cantarella.
\newblock A general mutual helicity formula.
\newblock {\em R. Soc. Lond. Proc. Ser. A Math. Phys. Eng. Sci.},
  456(2003):2771--2779, 2000.

\bibitem{Cantarella-DeTurck-Gluck01}
J.~Cantarella, D.~DeTurck, and H.~Gluck.
\newblock The {B}iot-{S}avart operator for application to knot theory, fluid
  dynamics, and plasma physics.
\newblock {\em J. Math. Phys.}, 42(2):876--905, 2001.

\bibitem{Cantarella-DeTurck-Gluck-Teytel00}
J.~Cantarella, D.~DeTurck, H.~Gluck, and M.~Teytel.
\newblock Isoperimetric problems for the helicity of vector fields and the
  {B}iot-{S}avart and curl operators.
\newblock {\em J. Math. Phys.}, 41(8):5615--5641, 2000.

\bibitem{Chavel84}
I.~Chavel.
\newblock {\em Eigenvalues in {R}iemannian geometry}, volume 115 of {\em Pure
  and Applied Mathematics}.
\newblock Academic Press Inc., Orlando, FL, 1984.
\newblock Including a chapter by Burton Randol, With an appendix by Jozef
  Dodziuk.

\bibitem{Cohen-private}
F.~R. Cohen.
\newblock Private communication.
\newblock 2007.

\bibitem{Cohen02}
F.~R. Cohen and S.~Gitler.
\newblock On loop spaces of configuration spaces.
\newblock {\em Trans. Amer. Math. Soc.}, 354(5):1705--1748 (electronic), 2002.

\bibitem{Deturck-Gluck-Melvin-Shonkwiler08}
D.~DeTurck, H.~Gluck, R.~Komendarczyk, P.~Melvin, C.~Shonkwiler, and
  D.~Vela-Vick.
\newblock Triple linking numbers, {H}opf invariants and {I}ntegral formulas for
  three-component links.
\newblock {\em in preparation}.

\bibitem{Berger-Evans92}
N.~W. Evans and M.~A. Berger.
\newblock A hierarchy of linking integrals.
\newblock In {\em Topological aspects of the dynamics of fluids and plasmas
  (Santa Barbara, CA, 1991)}, volume 218 of {\em NATO Adv. Sci. Inst. Ser. E
  Appl. Sci.}, pages 237--248. Kluwer Acad. Publ., Dordrecht, 1992.

\bibitem{Freedman91-2}
M.~Freedman and Z.~He.
\newblock Divergence-free fields: energy and asymptotic crossing number.
\newblock {\em Ann. of Math. (2)}, 134(1):189--229, 1991.

\bibitem{Gambaudo-Ghys97}
J.-M. Gambaudo and {\'E}.~Ghys.
\newblock Enlacements asymptotiques.
\newblock {\em Topology}, 36(6):1355--1379, 1997.

\bibitem{Khesin05}
B.~Khesin.
\newblock Topological fluid dynamics.
\newblock {\em Notices Amer. Math. Soc.}, 52(1):9--19, 2005.

\bibitem{Kohno02}
T.~Kohno.
\newblock Loop spaces of configuration spaces and finite type invariants.
\newblock In {\em Invariants of knots and 3-manifolds (Kyoto, 2001)}, volume~4
  of {\em Geom. Topol. Monogr.}, pages 143--160 (electronic). Geom. Topol.
  Publ., Coventry, 2002.

\bibitem{Laurence-Avellaneda93}
P.~Laurence and M.~Avellaneda.
\newblock A {M}offatt-{A}rnold formula for the mutual helicity of linked flux
  tubes.
\newblock {\em Geophys. Astrophys. Fluid Dynam.}, 69(1-4):243--256, 1993.

\bibitem{Laurence-Stredulinsky00b}
P.~Laurence and E.~Stredulinsky.
\newblock Asymptotic {M}assey products, induced currents and {B}orromean torus
  links.
\newblock {\em J. Math. Phys.}, 41(5):3170--3191, 2000.

\bibitem{Laurence-Stredulinsky00a}
P.~Laurence and E.~Stredulinsky.
\newblock A lower bound for the energy of magnetic fields supported in linked
  tori.
\newblock {\em C. R. Acad. Sci. Paris S\'er. I Math.}, 331(3):201--206, 2000.

\bibitem{Mayer03}
C.~Mayer.
\newblock Topological link invariants of magnetic fields.
\newblock {\em Ph.D. thesis}, 2003.

\bibitem{Milnor54}
J.~Milnor.
\newblock Link groups.
\newblock {\em Ann. of Math. (2)}, 59:177--195, 1954.

\bibitem{Milnor57}
J.~Milnor.
\newblock Isotopy of links.
\newblock In R.~Fox, editor, {\em Algebraic Geometry and Topology}, pages
  280--306. Princeton University Press, 1957.

\bibitem{Milnor97}
J.~Milnor.
\newblock {\em Topology from the Differentiable Viewpoint}, chapter~7.
\newblock Princeton Landmarks in Mathematics and Physics. Princeton University
  Press, Princeton, NJ, 1997.
\newblock Revised reprint of 1965 original.

\bibitem{Priest84}
E.~Priest.
\newblock {\em Solar Magnetohydrodynamics}.
\newblock D.Redidel Publishing Comp., 1984.

\bibitem{Schwarz95}
G.~Schwarz.
\newblock {\em Hodge decomposition---a method for solving boundary value
  problems}, volume 1607 of {\em Lecture Notes in Mathematics}.
\newblock Springer-Verlag, Berlin, 1995.

\bibitem{Hornig04}
H.~v.~Bodecker and G.~Hornig.
\newblock Link invariants of electromagnetic fields.
\newblock {\em Phys. Rev. Lett.}, 92(3):030406, 4, 2004.

\bibitem{Verjovsky94}
A.~Verjovsky and R.~F. Vila~Freyer.
\newblock The {J}ones-{W}itten invariant for flows on a {$3$}-dimensional
  manifold.
\newblock {\em Comm. Math. Phys.}, 163(1):73--88, 1994.

\bibitem{Vogel03}
T.~Vogel.
\newblock On the asymptotic linking number.
\newblock {\em Proc. Amer. Math. Soc.}, 131(7):2289--2297 (electronic), 2003.

\end{thebibliography}



\end{document}
